% Practical chances: the paper. The prose follows docs/writing.md. Every
% number, table, figure and board is written by `lab paper practical-chances`
% into generated/, from the model's committed results, and read back with
% \V{key} and \block{name}.
\documentclass[11pt]{article}
% The preamble every paper of the lab shares: % The preamble every paper of the lab shares: % The preamble every paper of the lab shares: \input{../../../papers/preamble}
% from models/<model>/paper/, \input{../preamble} from papers/<name>/.
% Boards come from chessboard and xskak, plots from pgfplots. A paper prints no
% number by hand: `lab paper` writes every number, table and figure under the
% paper's generated/ folder, and paper.tex says what they mean.
\usepackage[LSBC4,T1]{fontenc}
\usepackage[utf8]{inputenc}
\usepackage{lmodern}
\usepackage[margin=1in]{geometry}
\usepackage{microtype}
\usepackage{amsmath}
\usepackage{booktabs}
\usepackage{tabularx}
\usepackage{array}
\usepackage{xcolor}
\usepackage[obeyspaces,spaces]{url}
\usepackage{xskak}
\usepackage{pgfplots}
\usepgfplotslibrary{groupplots}
\usetikzlibrary{positioning,arrows.meta,calc,fit}
\pgfplotsset{compat=1.18}
\usepackage{caption}
\usepackage{subcaption}
\usepackage{enumitem}
\usepackage[section]{placeins}
\usepackage[hidelinks]{hyperref}

% A generated value by its key: generated/values.tex sets them, \V{key} prints
% one, and a key nobody set stops the build.
\makeatletter
\newcommand{\setvalue}[2]{\expandafter\gdef\csname labvalue@#1\endcsname{#2}}
\newcommand{\V}[1]{\@ifundefined{labvalue@#1}%
  {\PackageError{lab}{No generated value named '#1'}{Run `lab paper` for this model again.}}%
  {\csname labvalue@#1\endcsname}}
\makeatother

% Chart colours. A run keeps its colour in every figure of every paper.
\definecolor{labBlue}{HTML}{2A78D6}
\definecolor{labOrange}{HTML}{EB6834}
\definecolor{labAqua}{HTML}{1BAF7A}
\definecolor{labYellow}{HTML}{EDA100}
\definecolor{labMagenta}{HTML}{E87BA4}
\definecolor{labGreen}{HTML}{008300}
\definecolor{labViolet}{HTML}{6250D6}
\definecolor{labGray}{HTML}{898781}
\definecolor{labInk}{HTML}{0B0B0B}
\definecolor{labSecondary}{HTML}{52514E}
\definecolor{labGrid}{HTML}{E1E0D9}
\definecolor{labAxis}{HTML}{C3C2B7}
\definecolor{labWash}{HTML}{CDE2FB}

\pgfplotsset{
  lab/.style={
    font=\sffamily\footnotesize,
    axis line style={labAxis},
    tick style={labAxis},
    tick label style={labSecondary},
    label style={labSecondary},
    title style={labInk, font=\sffamily\small\bfseries, align=left},
    grid style={labGrid, line width=0.3pt},
    major grid style={labGrid, line width=0.3pt},
    axis lines*=left,
    tick align=outside,
    clip=false,
    legend style={draw=none, fill=none, font=\sffamily\footnotesize, text=labSecondary,
      cells={anchor=west}, legend columns=-1, column sep=6pt, at={(0,1.03)}, anchor=south west},
    every axis plot/.append style={line width=1.1pt, line join=round, line cap=round},
    scaled ticks=false,
    /pgf/number format/1000 sep={,},
  },
  lab marks/.style={mark=*, mark size=1.9pt, mark options={draw=white, line width=0.6pt}},
}

\setchessboard{boardfontencoding=LSBC4, boardfontfamily=skaknew, blackfieldcolor=black!18,
  whitefieldcolor=white, setfontcolors, showmover=false, label=true, labelfontsize=6pt,
  boardfontsize=12pt}

% A move or line in UCI, a position in FEN (breakable at a slash), and a run's name.
\newcommand{\uci}[1]{\texttt{#1}}
\newcommand{\fen}[1]{{\footnotesize\path{#1}}}
\newcommand{\run}[1]{\texttt{#1}}
\newcolumntype{L}{>{\raggedright\arraybackslash}X}
\newcolumntype{R}{>{\raggedleft\arraybackslash}X}

\setlength{\parskip}{0.4em}
\setlength{\parindent}{0pt}
\captionsetup{font=small, labelfont=bf}

% from models/<model>/paper/, \input{../preamble} from papers/<name>/.
% Boards come from chessboard and xskak, plots from pgfplots. A paper prints no
% number by hand: `lab paper` writes every number, table and figure under the
% paper's generated/ folder, and paper.tex says what they mean.
\usepackage[LSBC4,T1]{fontenc}
\usepackage[utf8]{inputenc}
\usepackage{lmodern}
\usepackage[margin=1in]{geometry}
\usepackage{microtype}
\usepackage{amsmath}
\usepackage{booktabs}
\usepackage{tabularx}
\usepackage{array}
\usepackage{xcolor}
\usepackage[obeyspaces,spaces]{url}
\usepackage{xskak}
\usepackage{pgfplots}
\usepgfplotslibrary{groupplots}
\usetikzlibrary{positioning,arrows.meta,calc,fit}
\pgfplotsset{compat=1.18}
\usepackage{caption}
\usepackage{subcaption}
\usepackage{enumitem}
\usepackage[section]{placeins}
\usepackage[hidelinks]{hyperref}

% A generated value by its key: generated/values.tex sets them, \V{key} prints
% one, and a key nobody set stops the build.
\makeatletter
\newcommand{\setvalue}[2]{\expandafter\gdef\csname labvalue@#1\endcsname{#2}}
\newcommand{\V}[1]{\@ifundefined{labvalue@#1}%
  {\PackageError{lab}{No generated value named '#1'}{Run `lab paper` for this model again.}}%
  {\csname labvalue@#1\endcsname}}
\makeatother

% Chart colours. A run keeps its colour in every figure of every paper.
\definecolor{labBlue}{HTML}{2A78D6}
\definecolor{labOrange}{HTML}{EB6834}
\definecolor{labAqua}{HTML}{1BAF7A}
\definecolor{labYellow}{HTML}{EDA100}
\definecolor{labMagenta}{HTML}{E87BA4}
\definecolor{labGreen}{HTML}{008300}
\definecolor{labViolet}{HTML}{6250D6}
\definecolor{labGray}{HTML}{898781}
\definecolor{labInk}{HTML}{0B0B0B}
\definecolor{labSecondary}{HTML}{52514E}
\definecolor{labGrid}{HTML}{E1E0D9}
\definecolor{labAxis}{HTML}{C3C2B7}
\definecolor{labWash}{HTML}{CDE2FB}

\pgfplotsset{
  lab/.style={
    font=\sffamily\footnotesize,
    axis line style={labAxis},
    tick style={labAxis},
    tick label style={labSecondary},
    label style={labSecondary},
    title style={labInk, font=\sffamily\small\bfseries, align=left},
    grid style={labGrid, line width=0.3pt},
    major grid style={labGrid, line width=0.3pt},
    axis lines*=left,
    tick align=outside,
    clip=false,
    legend style={draw=none, fill=none, font=\sffamily\footnotesize, text=labSecondary,
      cells={anchor=west}, legend columns=-1, column sep=6pt, at={(0,1.03)}, anchor=south west},
    every axis plot/.append style={line width=1.1pt, line join=round, line cap=round},
    scaled ticks=false,
    /pgf/number format/1000 sep={,},
  },
  lab marks/.style={mark=*, mark size=1.9pt, mark options={draw=white, line width=0.6pt}},
}

\setchessboard{boardfontencoding=LSBC4, boardfontfamily=skaknew, blackfieldcolor=black!18,
  whitefieldcolor=white, setfontcolors, showmover=false, label=true, labelfontsize=6pt,
  boardfontsize=12pt}

% A move or line in UCI, a position in FEN (breakable at a slash), and a run's name.
\newcommand{\uci}[1]{\texttt{#1}}
\newcommand{\fen}[1]{{\footnotesize\path{#1}}}
\newcommand{\run}[1]{\texttt{#1}}
\newcolumntype{L}{>{\raggedright\arraybackslash}X}
\newcolumntype{R}{>{\raggedleft\arraybackslash}X}

\setlength{\parskip}{0.4em}
\setlength{\parindent}{0pt}
\captionsetup{font=small, labelfont=bf}

% from models/<model>/paper/, \input{../preamble} from papers/<name>/.
% Boards come from chessboard and xskak, plots from pgfplots. A paper prints no
% number by hand: `lab paper` writes every number, table and figure under the
% paper's generated/ folder, and paper.tex says what they mean.
\usepackage[LSBC4,T1]{fontenc}
\usepackage[utf8]{inputenc}
\usepackage{lmodern}
\usepackage[margin=1in]{geometry}
\usepackage{microtype}
\usepackage{amsmath}
\usepackage{booktabs}
\usepackage{tabularx}
\usepackage{array}
\usepackage{xcolor}
\usepackage[obeyspaces,spaces]{url}
\usepackage{xskak}
\usepackage{pgfplots}
\usepgfplotslibrary{groupplots}
\usetikzlibrary{positioning,arrows.meta,calc,fit}
\pgfplotsset{compat=1.18}
\usepackage{caption}
\usepackage{subcaption}
\usepackage{enumitem}
\usepackage[section]{placeins}
\usepackage[hidelinks]{hyperref}

% A generated value by its key: generated/values.tex sets them, \V{key} prints
% one, and a key nobody set stops the build.
\makeatletter
\newcommand{\setvalue}[2]{\expandafter\gdef\csname labvalue@#1\endcsname{#2}}
\newcommand{\V}[1]{\@ifundefined{labvalue@#1}%
  {\PackageError{lab}{No generated value named '#1'}{Run `lab paper` for this model again.}}%
  {\csname labvalue@#1\endcsname}}
\makeatother

% Chart colours. A run keeps its colour in every figure of every paper.
\definecolor{labBlue}{HTML}{2A78D6}
\definecolor{labOrange}{HTML}{EB6834}
\definecolor{labAqua}{HTML}{1BAF7A}
\definecolor{labYellow}{HTML}{EDA100}
\definecolor{labMagenta}{HTML}{E87BA4}
\definecolor{labGreen}{HTML}{008300}
\definecolor{labViolet}{HTML}{6250D6}
\definecolor{labGray}{HTML}{898781}
\definecolor{labInk}{HTML}{0B0B0B}
\definecolor{labSecondary}{HTML}{52514E}
\definecolor{labGrid}{HTML}{E1E0D9}
\definecolor{labAxis}{HTML}{C3C2B7}
\definecolor{labWash}{HTML}{CDE2FB}

\pgfplotsset{
  lab/.style={
    font=\sffamily\footnotesize,
    axis line style={labAxis},
    tick style={labAxis},
    tick label style={labSecondary},
    label style={labSecondary},
    title style={labInk, font=\sffamily\small\bfseries, align=left},
    grid style={labGrid, line width=0.3pt},
    major grid style={labGrid, line width=0.3pt},
    axis lines*=left,
    tick align=outside,
    clip=false,
    legend style={draw=none, fill=none, font=\sffamily\footnotesize, text=labSecondary,
      cells={anchor=west}, legend columns=-1, column sep=6pt, at={(0,1.03)}, anchor=south west},
    every axis plot/.append style={line width=1.1pt, line join=round, line cap=round},
    scaled ticks=false,
    /pgf/number format/1000 sep={,},
  },
  lab marks/.style={mark=*, mark size=1.9pt, mark options={draw=white, line width=0.6pt}},
}

\setchessboard{boardfontencoding=LSBC4, boardfontfamily=skaknew, blackfieldcolor=black!18,
  whitefieldcolor=white, setfontcolors, showmover=false, label=true, labelfontsize=6pt,
  boardfontsize=12pt}

% A move or line in UCI, a position in FEN (breakable at a slash), and a run's name.
\newcommand{\uci}[1]{\texttt{#1}}
\newcommand{\fen}[1]{{\footnotesize\path{#1}}}
\newcommand{\run}[1]{\texttt{#1}}
\newcolumntype{L}{>{\raggedright\arraybackslash}X}
\newcolumntype{R}{>{\raggedleft\arraybackslash}X}

\setlength{\parskip}{0.4em}
\setlength{\parindent}{0pt}
\captionsetup{font=small, labelfont=bf}

% Written by `lab paper practical-chances`; do not edit by hand.
\looktrue
\chancefalse
\sequencefalse
\contextgbdttrue
\pairstrue
\variantstrue
\shapesfalse
\gamestrue
\featurestrue
\changestrue
\changehandtrue
\changeepochsfalse
\changestarttrue
\changeattacksfalse
\changegeometryfalse
\changeconvfalse
\changehybridfalse
\changelargerfalse
\moretrue
\shownfalse
\shownpairedfalse
\setvalue{band.2.floor}{0.4949}
\setvalue{band.4.floor}{0.4917}
\setvalue{band.5.floor}{0.4809}
\setvalue{band.5.name}{2000+}
\setvalue{bands.best-ahead.of}{all five bands}
\setvalue{bands.count}{five}
\setvalue{bands.lead.narrowest}{+0.000149}
\setvalue{bands.lead.narrowest.name}{1700-2000}
\setvalue{bands.lead.widest}{+0.00315}
\setvalue{bands.lead.widest.name}{2000+}
\setvalue{best.beside.correlation}{0.99}
\setvalue{best.gain}{+0.001353}
\setvalue{best.gain.high}{+0.002213}
\setvalue{best.gain.low}{+0.000469}
\setvalue{best.gain.of-lift}{2}
\setvalue{best.gain.percent}{0.32}
\setvalue{best.name}{\run{gbm-hand-quadruple}}
\setvalue{best.r2}{0.2619}
\setvalue{best.reliability.groups}{20}
\setvalue{best.reliability.highest.happened}{0.966}
\setvalue{best.reliability.highest.said}{0.967}
\setvalue{best.reliability.lowest.happened}{0.030}
\setvalue{best.reliability.lowest.said}{0.030}
\setvalue{best.reliability.widest-gap}{0.030}
\setvalue{best.score}{0.4223}
\setvalue{best.score.high}{0.4281}
\setvalue{best.score.low}{0.4169}
\setvalue{best.seconds}{42 s}
\setvalue{best.spearman}{0.5063}
\setvalue{best.trees}{327}
\setvalue{challenger.name}{\run{gbm-hand-quadruple}}
\setvalue{code.commit}{\texttt{173ddde}}
\setvalue{context.lift.share}{97}
\setvalue{cost.seconds}{33 min}
\setvalue{cost.seconds.networks}{32 min}
\setvalue{cost.seconds.others}{69 s}
\setvalue{data.fingerprint}{\texttt{340c68ed5c91}}
\setvalue{data.games.test}{2{,}944}
\setvalue{data.licence}{CC0 1.0}
\setvalue{data.month}{October 2026}
\setvalue{data.name}{\texttt{game-positions-2026-10}}
\setvalue{data.rows}{698{,}013}
\setvalue{data.rows.test}{68{,}579}
\setvalue{data.share.test}{10}
\setvalue{data.share.train}{85}
\setvalue{data.share.valid}{5}
\setvalue{describe.clock.q1}{28}
\setvalue{describe.clock.q5}{139}
\setvalue{describe.clock.q50}{506}
\setvalue{describe.loss.mean}{4.71}
\setvalue{describe.win-before.max}{97.5}
\setvalue{describe.win-before.mean}{52.4}
\setvalue{describe.win-before.min}{2.46}
\setvalue{describe.win-before.q25}{31.0}
\setvalue{describe.win-before.q75}{79.5}
\setvalue{examples.count}{three}
\setvalue{floor.constant}{0.5013}
\setvalue{floor.score}{0.4916}
\setvalue{grid.clock-win.-30.40-60.happened}{0.324}
\setvalue{grid.clock-win.-30.40-60.rows}{139}
\setvalue{grid.clock-win.-30.80+.engine-curve.said}{0.947}
\setvalue{grid.clock-win.-30.80+.happened}{0.723}
\setvalue{grid.clock-win.30-60.60-80.rows}{50}
\setvalue{grid.clock-win.300+.80+.engine-curve.said}{0.918}
\setvalue{grid.clock-win.300+.80+.happened}{0.805}
\setvalue{grid.clock-win.engine-curve.widest-off}{0.31}
\setvalue{grid.clock-win.engine-curve.widest-off.column}{60-80}
\setvalue{grid.clock-win.engine-curve.widest-off.row}{30-60}
\setvalue{grid.rating-win.-1100.-20.engine-curve.said}{0.070}
\setvalue{grid.rating-win.-1100.-20.happened}{0.193}
\setvalue{grid.rating-win.-1100.60-80.happened}{0.586}
\setvalue{grid.rating-win.-1100.80+.engine-curve.said}{0.932}
\setvalue{grid.rating-win.-1100.80+.happened}{0.760}
\setvalue{grid.rating-win.2000+.-20.engine-curve.said}{0.093}
\setvalue{grid.rating-win.2000+.-20.happened}{0.090}
\setvalue{grid.rating-win.2000+.60-80.happened}{0.722}
\setvalue{grid.rating-win.2000+.80+.engine-curve.said}{0.912}
\setvalue{grid.rating-win.2000+.80+.happened}{0.901}
\setvalue{grid.rating-win.engine-curve.widest-off}{0.17}
\setvalue{grid.rating-win.engine-curve.widest-off.column}{80+}
\setvalue{grid.rating-win.engine-curve.widest-off.row}{\textless{}1100}
\setvalue{grid.rating-win.gbm-hand.widest-off}{0.041}
\setvalue{grid.smallest-cell}{30}
\setvalue{label.mean}{0.5013}
\setvalue{label.sd}{0.4910}
\setvalue{ladder.at-last-epoch.of}{none of the four networks}
\setvalue{ladder.count}{four}
\setvalue{look.date}{23 November 2026}
\setvalue{look.resamples}{200}
\setvalue{look.smallest-compared}{200}
\setvalue{look.test-rows}{68{,}579}
\setvalue{manifest.games-read}{1{,}003{,}668}
\setvalue{method.gbm.features.percent}{80}
\setvalue{method.gbm.leaves}{127}
\setvalue{method.gbm.patience}{100}
\setvalue{method.gbm.rate}{0.05}
\setvalue{method.gbm.rounds}{4{,}000}
\setvalue{method.gbm.rows.percent}{80}
\setvalue{method.resamples}{1{,}000}
\setvalue{more.gbm-hand.most}{four times}
\setvalue{more.transformer-medium-hand.most}{four times}
\setvalue{network.device}{NVIDIA A100-SXM4-40GB}
\setvalue{network.epochs.run}{3}
\setvalue{network.gain}{+0.001210}
\setvalue{network.gain.high}{+0.002466}
\setvalue{network.gain.low}{$-$0.000025}
\setvalue{network.name}{\run{transformer-medium-hand-quadruple}}
\setvalue{network.parameters}{6{,}557{,}505}
\setvalue{network.rows-a-second}{14{,}000}
\setvalue{network.score}{0.4225}
\setvalue{network.seconds}{9 min}
\setvalue{networks.ahead.of}{none of the six networks}
\setvalue{networks.behind.of}{none of the six networks}
\setvalue{networks.patience}{three}
\setvalue{networks.warmup.percent}{5}
\setvalue{overall.engine-curve.off}{+0.023}
\setvalue{overall.engine-curve.said}{0.525}
\setvalue{overall.gbm-hand.off}{$-$0.00023}
\setvalue{overall.happened}{0.502}
\setvalue{pair.gbm-context.elo-gap.gain}{+0.07069}
\setvalue{pair.gbm-context.elo-gap.gain.high}{+0.07692}
\setvalue{pair.gbm-context.elo-gap.gain.low}{+0.06454}
\setvalue{pair.gbm-context.elo-gap.verdict}{significantly better}
\setvalue{pair.gbm-context.engine-curve.gain}{+0.01136}
\setvalue{pair.gbm-context.engine-curve.gain.high}{+0.01369}
\setvalue{pair.gbm-context.engine-curve.gain.low}{+0.00892}
\setvalue{pair.gbm-context.engine-curve.gain.percent}{2.6}
\setvalue{pair.gbm-context.engine-curve.verdict}{significantly better}
\setvalue{pair.gbm-hand-double.gbm-hand.gain}{+0.000521}
\setvalue{pair.gbm-hand-double.gbm-hand.kept}{not accepted}
\setvalue{pair.gbm-hand-double.gbm-hand.verdict}{not significantly different}
\setvalue{pair.gbm-hand-half.gbm-hand.gain}{$-$0.000475}
\setvalue{pair.gbm-hand-half.gbm-hand.verdict}{not significantly different}
\setvalue{pair.gbm-hand-no-hanging.gbm-hand.gain}{$-$0.0002129}
\setvalue{pair.gbm-hand-no-hanging.gbm-hand.verdict}{not significantly different}
\setvalue{pair.gbm-hand-no-kings.gbm-hand.gain}{$-$0.0001212}
\setvalue{pair.gbm-hand-no-kings.gbm-hand.verdict}{not significantly different}
\setvalue{pair.gbm-hand-no-mobility.gbm-hand.gain}{$-$0.000550}
\setvalue{pair.gbm-hand-no-mobility.gbm-hand.verdict}{significantly worse}
\setvalue{pair.gbm-hand-no-pawns.gbm-hand.gain}{$-$0.0001402}
\setvalue{pair.gbm-hand-no-pawns.gbm-hand.verdict}{not significantly different}
\setvalue{pair.gbm-hand-quadruple.gbm-hand.gain}{+0.001353}
\setvalue{pair.gbm-hand-quadruple.gbm-hand.kept}{accepted}
\setvalue{pair.gbm-hand-quadruple.gbm-hand.verdict}{significantly better}
\setvalue{pair.gbm-hand-quarter.gbm-hand.gain}{$-$0.000681}
\setvalue{pair.gbm-hand-quarter.gbm-hand.verdict}{not significantly different}
\setvalue{pair.gbm-hand.gbm-context.gain}{+0.002288}
\setvalue{pair.gbm-hand.gbm-context.gain.high}{+0.003088}
\setvalue{pair.gbm-hand.gbm-context.gain.low}{+0.001502}
\setvalue{pair.gbm-hand.gbm-context.kept}{accepted}
\setvalue{pair.gbm-hand.gbm-context.valid.verdict}{better}
\setvalue{pair.gbm-hand.gbm-context.verdict}{significantly better}
\setvalue{pair.transformer-medium-hand-double.transformer-medium-hand.gain}{+0.001403}
\setvalue{pair.transformer-medium-hand-double.transformer-medium-hand.kept}{accepted}
\setvalue{pair.transformer-medium-hand-double.transformer-medium-hand.verdict}{significantly better}
\setvalue{pair.transformer-medium-hand-half.transformer-medium-hand.gain}{$-$0.001023}
\setvalue{pair.transformer-medium-hand-pretrained.transformer-medium-hand.gain}{$-$0.001252}
\setvalue{pair.transformer-medium-hand-pretrained.transformer-medium-hand.gain.high}{$-$0.000165}
\setvalue{pair.transformer-medium-hand-pretrained.transformer-medium-hand.gain.low}{$-$0.002356}
\setvalue{pair.transformer-medium-hand-pretrained.transformer-medium-hand.kept}{not accepted}
\setvalue{pair.transformer-medium-hand-pretrained.transformer-medium-hand.valid.verdict}{no better}
\setvalue{pair.transformer-medium-hand-pretrained.transformer-medium-hand.verdict}{significantly worse}
\setvalue{pair.transformer-medium-hand-quadruple.transformer-medium-hand.gain}{+0.002430}
\setvalue{pair.transformer-medium-hand-quadruple.transformer-medium-hand.kept}{accepted}
\setvalue{pair.transformer-medium-hand-quadruple.transformer-medium-hand.verdict}{significantly better}
\setvalue{pair.transformer-medium-hand-quarter.transformer-medium-hand.gain}{$-$0.002697}
\setvalue{pair.transformer-medium-hand.engine-curve.gain}{+0.01243}
\setvalue{pair.transformer-medium-hand.engine-curve.gain.high}{+0.01502}
\setvalue{pair.transformer-medium-hand.engine-curve.gain.low}{+0.00992}
\setvalue{pair.transformer-medium-hand.engine-curve.verdict}{significantly better}
\setvalue{pair.transformer-medium-hand.transformer-medium.gain}{$-$0.000479}
\setvalue{pair.transformer-medium-hand.transformer-medium.gain.high}{+0.000368}
\setvalue{pair.transformer-medium-hand.transformer-medium.gain.low}{$-$0.001288}
\setvalue{pair.transformer-medium-hand.transformer-medium.kept}{not accepted}
\setvalue{pair.transformer-medium-hand.transformer-medium.valid.verdict}{no better}
\setvalue{pair.transformer-medium-hand.transformer-medium.verdict}{not significantly different}
\setvalue{pair.transformer-medium-pretrained.transformer-medium.gain}{$-$0.002151}
\setvalue{pair.transformer-medium-pretrained.transformer-medium.gain.high}{$-$0.000856}
\setvalue{pair.transformer-medium-pretrained.transformer-medium.gain.low}{$-$0.003404}
\setvalue{pair.transformer-medium-pretrained.transformer-medium.kept}{not accepted}
\setvalue{pair.transformer-medium-pretrained.transformer-medium.valid.verdict}{no better}
\setvalue{pair.transformer-medium-pretrained.transformer-medium.verdict}{significantly worse}
\setvalue{recipe.games}{30{,}000}
\setvalue{recipe.min-ply}{11}
\setvalue{recipe.per-game}{24}
\setvalue{reference.features}{25}
\setvalue{reference.leans.1.share}{73}
\setvalue{reference.leans.2.feature}{opponent's rating}
\setvalue{reference.leans.2.share}{5}
\setvalue{reference.leans.3.feature}{mover's rating}
\setvalue{reference.leans.3.share}{5}
\setvalue{reference.leans.4.feature}{material balance}
\setvalue{reference.leans.4.share}{4}
\setvalue{reference.leans.5.feature}{opponent's clock}
\setvalue{reference.leans.5.share}{3}
\setvalue{reference.leans.6.feature}{mover's clock}
\setvalue{reference.leans.6.share}{3}
\setvalue{reference.mae}{0.3618}
\setvalue{reference.r2}{0.2572}
\setvalue{reference.reliability.groups}{20}
\setvalue{reference.reliability.widest-gap}{0.026}
\setvalue{reference.score}{0.4237}
\setvalue{reference.seconds}{22 s}
\setvalue{reference.spearman}{0.5016}
\setvalue{reference.trees}{208}
\setvalue{results.date}{23 November 2026}
\setvalue{run.conv-small.best-epoch}{5}
\setvalue{run.conv-small.epochs}{8}
\setvalue{run.conv-small.gain}{+0.000162}
\setvalue{run.conv-small.gain.high}{+0.001411}
\setvalue{run.conv-small.gain.low}{$-$0.001151}
\setvalue{run.elo-gap.band.5}{0.4655}
\setvalue{run.elo-gap.beside.correlation}{0.16}
\setvalue{run.elo-gap.r2}{$-$0.0207}
\setvalue{run.elo-gap.r2.high}{$-$0.0030}
\setvalue{run.elo-gap.r2.low}{$-$0.0406}
\setvalue{run.elo-gap.reliability.highest.happened}{0.603}
\setvalue{run.elo-gap.reliability.highest.said}{0.809}
\setvalue{run.elo-gap.reliability.lowest.happened}{0.417}
\setvalue{run.elo-gap.reliability.lowest.said}{0.190}
\setvalue{run.elo-gap.score}{0.4967}
\setvalue{run.elo-gap.score.high}{0.5016}
\setvalue{run.elo-gap.score.low}{0.4918}
\setvalue{run.elo-gap.spearman}{0.0974}
\setvalue{run.engine-curve.beside.closer}{58}
\setvalue{run.engine-curve.beside.correlation}{0.95}
\setvalue{run.engine-curve.gain}{$-$0.01365}
\setvalue{run.engine-curve.gain.high}{$-$0.01124}
\setvalue{run.engine-curve.gain.low}{$-$0.01612}
\setvalue{run.engine-curve.lift.share}{80}
\setvalue{run.engine-curve.mae}{0.3487}
\setvalue{run.engine-curve.r2}{0.2086}
\setvalue{run.engine-curve.reliability.highest.happened}{0.924}
\setvalue{run.engine-curve.reliability.highest.said}{0.975}
\setvalue{run.engine-curve.reliability.lowest.happened}{0.055}
\setvalue{run.engine-curve.reliability.lowest.said}{0.025}
\setvalue{run.engine-curve.reliability.widest-gap}{0.16}
\setvalue{run.engine-curve.reliability.widest-gap.happened}{0.670}
\setvalue{run.engine-curve.reliability.widest-gap.said}{0.834}
\setvalue{run.engine-curve.score}{0.4373}
\setvalue{run.engine-curve.score.high}{0.4439}
\setvalue{run.engine-curve.score.low}{0.4313}
\setvalue{run.engine-curve.spearman}{0.4780}
\setvalue{run.gbm-context.gain}{$-$0.002288}
\setvalue{run.gbm-context.r2}{0.2492}
\setvalue{run.gbm-context.score}{0.4260}
\setvalue{run.gbm-context.spearman}{0.4936}
\setvalue{run.gbm-hand-quadruple.score}{0.4223}
\setvalue{run.gbm-hand-quadruple.trees}{327}
\setvalue{run.transformer-medium-hand-double.best-epoch}{3}
\setvalue{run.transformer-medium-hand-double.epochs}{6}
\setvalue{run.transformer-medium-hand-pretrained.best-epoch}{4}
\setvalue{run.transformer-medium-hand-quadruple.best-epoch}{3}
\setvalue{run.transformer-medium-hand-quadruple.epochs}{3}
\setvalue{run.transformer-medium-hand-quadruple.gain}{+0.001210}
\setvalue{run.transformer-medium-hand-quadruple.gain.high}{+0.002466}
\setvalue{run.transformer-medium-hand-quadruple.gain.low}{$-$0.000025}
\setvalue{run.transformer-medium-hand-quadruple.score}{0.4225}
\setvalue{run.transformer-medium-hand.best-epoch}{5}
\setvalue{run.transformer-medium-hand.epochs}{12}
\setvalue{run.transformer-medium-hand.gain}{$-$0.001220}
\setvalue{run.transformer-medium-hand.gain.high}{+0.000007}
\setvalue{run.transformer-medium-hand.gain.low}{$-$0.002381}
\setvalue{run.transformer-medium-pretrained.best-epoch}{3}
\setvalue{run.transformer-medium-pretrained.epochs}{12}
\setvalue{run.transformer-medium.best-epoch}{5}
\setvalue{run.transformer-medium.epochs}{12}
\setvalue{run.transformer-medium.epochs.run}{8}
\setvalue{run.transformer-medium.gain}{$-$0.000741}
\setvalue{run.transformer-medium.gain.high}{+0.000580}
\setvalue{run.transformer-medium.gain.low}{$-$0.002056}
\setvalue{run.transformer-small.best-epoch}{6}
\setvalue{run.transformer-small.epochs}{8}
\setvalue{run.transformer-small.gain}{+0.000157}
\setvalue{run.transformer-small.gain.high}{+0.001565}
\setvalue{run.transformer-small.gain.low}{$-$0.001199}
\setvalue{runs.count}{twelve}
\setvalue{slice.band.ahead.of}{two of the five groups}
\setvalue{slice.band.engine-curve.behind.least}{+0.00641}
\setvalue{slice.band.engine-curve.behind.least.group}{1700-2000}
\setvalue{slice.band.engine-curve.behind.most}{+0.0255}
\setvalue{slice.band.engine-curve.behind.most.group}{\textless{}1100}
\setvalue{slice.clock.-30.engine-curve.said}{0.480}
\setvalue{slice.clock.-30.happened}{0.383}
\setvalue{slice.clock.-30.reference.said}{0.358}
\setvalue{slice.clock.-30.rows}{718}
\setvalue{slice.clock.120-300.happened}{0.509}
\setvalue{slice.clock.30-60.engine-curve.behind}{+0.0631}
\setvalue{slice.clock.30-60.engine-curve.said}{0.515}
\setvalue{slice.clock.30-60.happened}{0.388}
\setvalue{slice.clock.30-60.rows}{777}
\setvalue{slice.clock.300-480.happened}{0.506}
\setvalue{slice.clock.480+.engine-curve.behind}{+0.00949}
\setvalue{slice.clock.480+.happened}{0.505}
\setvalue{slice.clock.60-120.engine-curve.behind}{+0.0486}
\setvalue{slice.clock.60-120.happened}{0.453}
\setvalue{slice.clock.60-120.rows}{1{,}456}
\setvalue{slice.clock.reference.widest-off}{0.033}
\setvalue{slice.clock.reference.widest-off.group}{30-60}
\setvalue{slice.hanging.0.engine-curve.said}{0.536}
\setvalue{slice.hanging.0.happened}{0.514}
\setvalue{slice.hanging.2.engine-curve.said}{0.480}
\setvalue{slice.hanging.2.happened}{0.462}
\setvalue{slice.opponent-clock.-30.engine-curve.said}{0.586}
\setvalue{slice.opponent-clock.-30.happened}{0.666}
\setvalue{slice.opponent-clock.-30.reference.said}{0.680}
\setvalue{slice.opponent-hanging.2.engine-curve.said}{0.583}
\setvalue{slice.opponent-hanging.2.happened}{0.554}
\setvalue{slice.ply.-20.constant}{0.4952}
\setvalue{slice.ply.-20.engine-curve.behind}{+0.00880}
\setvalue{slice.ply.-20.reference}{0.4775}
\setvalue{slice.ply.80+.constant}{0.4655}
\setvalue{slice.ply.80+.engine-curve.behind}{+0.0268}
\setvalue{slice.ply.80+.reference}{0.2827}
\setvalue{slice.win.-10.engine-curve.said}{0.050}
\setvalue{slice.win.-10.happened}{0.109}
\setvalue{slice.win.10-20.challenger.off}{$-$0.036}
\setvalue{slice.win.10-20.engine-curve.said}{0.144}
\setvalue{slice.win.10-20.happened}{0.266}
\setvalue{slice.win.40-50.engine-curve.off}{+0.037}
\setvalue{slice.win.50-60.engine-curve.off}{+0.0078}
\setvalue{slice.win.70-80.engine-curve.said}{0.751}
\setvalue{slice.win.80-90.challenger.off}{+0.032}
\setvalue{slice.win.80-90.engine-curve.said}{0.856}
\setvalue{slice.win.80-90.happened}{0.699}
\setvalue{slice.win.80-90.reference.off}{+0.033}
\setvalue{slice.win.90+.engine-curve.said}{0.951}
\setvalue{slice.win.90+.happened}{0.859}
\setvalue{slice.win.engine-curve.behind.least}{+0.00468}
\setvalue{slice.win.engine-curve.behind.least.group}{50-60}
\setvalue{slice.win.engine-curve.behind.most}{+0.0348}
\setvalue{slice.win.engine-curve.behind.most.group}{80-90}
\setvalue{slice.win.engine-curve.widest-off}{0.16}
\setvalue{slice.win.engine-curve.widest-off.group}{80-90}
\setvalue{slice.win.reference.widest-off}{0.035}
\setvalue{slice.win.reference.widest-off.group}{10-20}
\setvalue{slices.ahead}{27}
\setvalue{slices.behind}{zero}
\setvalue{slices.compared}{48}
\setvalue{study.features.gbm-hand.most.gain}{$-$0.001706}
\setvalue{study.features.gbm-hand.most.group}{material}
\setvalue{study.features.gbm-hand.worse.of}{two of the five ablations}
\setvalue{study.games.transformer-medium-hand.worse.of}{both ablations}
\setvalue{task.context.count}{twelve}
\setvalue{task.hand.count}{13}
\setvalue{task.marks.count}{three}
\setvalue{task.reference}{\run{gbm-hand}}


\begin{document}

\papertitle{Predicting the Results of Human Chess Games from Engine
  Evaluations, Ratings and Clocks}
  {Results as of \V{results.date}\quad\textperiodcentered\quad dataset \V{data.name},
   build \V{data.fingerprint}\quad\textperiodcentered\quad code \V{code.commit}}

\begin{abstract}
A chess engine evaluates a position under best play by both sides, and chess
software commonly converts the evaluation into a win probability. We measure
how well this win probability predicts the results of games between human
players, and we train models that predict the expected score of the player to
move, where a win counts 1, a draw one half and a loss 0, from the position,
the engine-derived win probability, both players' ratings, both clocks and
the time control. The dataset contains \V{data.rows} positions from
\V{recipe.games} rated rapid games played on Lichess, split by game, and the
primary metric is the root mean squared error (RMSE) of the predicted score.
Used directly as a prediction, the win probability has an RMSE of
\V{run.engine-curve.score} on \V{data.rows.test} test positions, compared
with \V{floor.score} for a constant prediction. Gradient-boosted decision
trees (GBDT) on the win probability, the ratings, the clocks and the ply
reach \V{run.gbm-context.score}, a paired gain of
\V{pair.gbm-context.engine-curve.gain} (95\% CI
\V{pair.gbm-context.engine-curve.gain.low} to
\V{pair.gbm-context.engine-curve.gain.high}), and GBDT that also receives
hand-crafted features of the position reaches \V{reference.score}. None of
the convolutional and transformer networks that read the board is
significantly better than this GBDT baseline, and pretraining a transformer's
encoder on engine evaluations makes it significantly worse. The best of
\V{runs.count} models is the GBDT baseline trained on \V{more.gbm-hand.most}
as many games, with an RMSE of \V{best.score} (95\% CI \V{best.score.low} to
\V{best.score.high}).\iflook{} The engine-derived win probability is
overconfident at both ends of its range, most for lower-rated players: from
positions that it rates at 80\% or more, movers rated below 1100 score
\V{grid.rating-win.-1100.80+.happened} and movers rated 2000 and above score
\V{grid.rating-win.2000+.80+.happened}. It also takes no account of the
clock: with less than 30 seconds remaining, movers score
\V{slice.clock.-30.happened}, where it gives them
\V{slice.clock.-30.engine-curve.said} on average.\fi{} The Elo formula applied
to the two ratings alone predicts the results less accurately than a
constant. The coefficient of determination of the best model is \V{best.r2}:
a single position, even with its evaluation, leaves most of the variation in
the result of a game unexplained.
\end{abstract}

\section{Introduction}
\label{sec:intro}

A chess engine assigns every position an evaluation: the value of the
position if both sides play the best moves from then on. Analysis software
converts this evaluation into a win probability with a fixed curve and
displays it next to the position~\cite{lichessaccuracy}. A player who reads
``85\%'' takes it as a statement about the game: I should have won this.
Whether that reading is justified depends on who is playing. A winning
position must still be won and a lost one can still be saved, and how often
this happens depends on the players and on their clocks.

This work supports a chess coaching application that reviews a player's
games. Its feedback refers to how the player stood, and for that it needs the
result that players of a given rating actually obtain from a position, next
to the result under best play. Where the two agree, the engine's verdict can
be reported as it is. Where they differ, the difference is informative: a
position that the engine considers won and that players of the user's rating
often fail to win is a position worth studying.

We study the following prediction task. Before a move is played, given the
position, the engine-derived win probability of the player to move (the
\emph{mover}), both players' ratings, both clock times and the time control,
predict the mover's expected score in the game: the mean result, with a win
counted as 1, a draw as one half and a loss as 0. The win probability used
directly as a prediction is one baseline, and the Elo expected score computed
from the two ratings is another. Neither has a fitted parameter.

We evaluate a fixed sequence of models on \V{data.rows} positions from rated
rapid games on Lichess: a constant baseline, the two baselines without fitted
parameters, GBDT on scalar context features, which include the win
probability, GBDT with hand-crafted features of the position, and neural
networks that read the board directly, with and without the hand-crafted
features. Our main findings are the following.
\begin{itemize}
\item \textbf{A fitted model is significantly more accurate than the
engine-derived win probability.} Used as a prediction, the win probability
has an RMSE of \V{run.engine-curve.score}. GBDT on the win probability, the
ratings, the clocks, the time control and the ply, without the board, reaches
\V{run.gbm-context.score}, a paired gain of
\V{pair.gbm-context.engine-curve.gain} (95\% CI
\V{pair.gbm-context.engine-curve.gain.low} to
\V{pair.gbm-context.engine-curve.gain.high}), which is
\V{pair.gbm-context.engine-curve.gain.percent}\% of the RMSE of the win
probability (Section~\ref{sec:results}).
\iflook
\item \textbf{The win probability is overconfident, most for lower-rated
players and under time pressure.} Where it gives the mover 80\% to 90\%,
movers score \V{slice.win.80-90.happened}, and where it gives 10\% to 20\%
they score \V{slice.win.10-20.happened}. From positions that it rates at 80\%
or more, movers rated below 1100 score \V{grid.rating-win.-1100.80+.happened}
and movers rated 2000 and above score \V{grid.rating-win.2000+.80+.happened}.
With less than 30 seconds on the mover's clock, movers score
\V{slice.clock.-30.happened}, where the win probability averages
\V{slice.clock.-30.engine-curve.said}. The mean prediction of the GBDT
baseline is within \V{reference.reliability.widest-gap} of the mean result
over the whole range of its predictions (Section~\ref{sec:error}).
\fi
\item \textbf{The position adds little to its evaluation.} The hand-crafted
features improve GBDT by a paired gain of \V{pair.gbm-hand.gbm-context.gain}
(95\% CI \V{pair.gbm-hand.gbm-context.gain.low} to
\V{pair.gbm-hand.gbm-context.gain.high}), and none of the \V{ladder.count}
networks of the fixed sequence differs significantly from the GBDT baseline.
GBDT on the context features alone achieves \V{context.lift.share}\% of the
GBDT baseline's improvement over the constant baseline, and the win
probability alone achieves \V{run.engine-curve.lift.share}\%.
\ifchangestart
\item \textbf{Pretraining on engine evaluations does not help.} A transformer
whose encoder is initialised from a model trained on engine evaluations is
significantly worse than the same transformer trained from random
initialisation, with and without the hand-crafted features
(Section~\ref{sec:ablations}).
\fi
\ifmore
\item \textbf{More training data improves both model classes.} With
\V{more.gbm-hand.most} the training games, the paired gain is
\V{pair.gbm-hand-quadruple.gbm-hand.gain} for GBDT and
\V{pair.transformer-medium-hand-quadruple.transformer-medium-hand.gain} for
the medium transformer, and the two models then have almost the same RMSE
(\V{run.gbm-hand-quadruple.score} and
\V{run.transformer-medium-hand-quadruple.score};
Section~\ref{sec:training-size}).
\fi
\item \textbf{The two ratings alone do not predict the results of these
games.} The Elo expected score computed from the two ratings has an RMSE of
\V{run.elo-gap.score}, higher than the \V{floor.score} of the constant
baseline, although the fitted models use both ratings
(Section~\ref{sec:rating-baseline}).
\item \textbf{Most of the variation is not predictable from one position.}
The coefficient of determination (R\textsuperscript{2}) of the best model is
\V{best.r2}. A prediction describes the mean result of many comparable games
and says little about a single game.
\end{itemize}

\section{Related Work}
\label{sec:related}

\paragraph{Engine evaluations as win probabilities.}
Lichess converts an engine evaluation into a win probability with a logistic
curve~\cite{lichessaccuracy}, whose single coefficient was fitted to
positions from rapid games between players rated 2300 and
above~\cite{lilawinchances}. The curve depends on the evaluation alone. We
measure how well it predicts the results of games across the whole range of
ratings, and how much the ratings, the clocks and the position add to it.

\paragraph{Predicting the outcome of a game.}
AlphaZero~\cite{alphazero} trains a network with a value head that predicts
the outcome of its own games. Allie~\cite{allie}, a transformer trained on
sequences of moves from human games, has a value head trained on the outcomes
of those games, and its authors report that the outcome becomes more
predictable as a game progresses. Maia~\cite{maia} and Maia-2~\cite{maia2}
predict the move that a player of a given rating plays. Our setting differs
in that every model is given an engine evaluation of the position, and the
question is what the ratings, the clocks and the board add to it.

\paragraph{Rating systems.}
The Elo formula gives the expected score of a game from the difference
between two ratings. Lichess rates its players with
Glicko-2~\cite{glicko2}, which keeps a rating deviation next to each rating
as a measure of its uncertainty. We apply the Elo formula to the two Lichess
ratings as a baseline that ignores the position.

\paragraph{Human error and time.}
Anderson et al.~\cite{humanerror} relate the errors of human players to the
difficulty of the position, the player's skill and the time available.
Companion papers of this repository model the risk of a blunder and the time
spent on a move for the positions used here. The present task concerns the
sum of all later errors of both players, as it appears in the result of the
game.

\paragraph{Pretraining on engine evaluations.}
Ruoss et al.~\cite{searchless} show that a large transformer trained by
supervised learning on engine evaluations plays strongly without search. We
use a far smaller model of this kind only as a source of pretrained weights
(Section~\ref{sec:ablations}).

\section{Task and Data}
\label{sec:data}

\subsection{Task}

An example is one position of a game, immediately before a move is played.
The input is the position, the opponent's preceding move, both players'
ratings, both clock times, the time control, the ply and the mover's
engine-derived win probability. The target is the result of the game for the
mover. Table~\ref{tab:brief} summarises the task.

\begin{table}[htbp]\centering\small
\begin{tabularx}{\textwidth}{lL}
\toprule
Input & The position before the move, the opponent's preceding move, both
players' ratings, both clock times, the time control, the ply and the mover's
win probability before the move. \\
Target & The result of the game for the mover: 1 for a win, 0.5 for a draw
and 0 for a loss. A model predicts its expected value, the expected score. \\
Data & Rated rapid games from the Lichess database (\V{data.licence}) for
\V{data.month} that carry engine evaluations and clock times:
\V{recipe.games} games and up to \V{recipe.per-game} positions per game. \\
Split & By game: \V{data.share.train}\% training, \V{data.share.valid}\%
validation, \V{data.share.test}\% test. \\
Primary metric & RMSE of the predicted score on the test set (lower is
better). \\
\bottomrule
\end{tabularx}
\caption{Summary of the task.}
\label{tab:brief}
\end{table}

\paragraph{Definition of the target.}
Let $r$ be the result of a game for White: 1 for a win, $\tfrac{1}{2}$ for a
draw and 0 for a loss. The target of a position is $y = r$ when White is to
move and $y = 1 - r$ when Black is to move. A game that was decided on time
has the result that was recorded for it. The target does not describe the
position. It is a single outcome of everything that both players did
afterwards, and all positions of a game share it up to the side to move. What
a model can estimate from an input $x$ is the expected score
$\mathrm{E}[y \mid x]$, the mean result of movers in comparable situations.

\paragraph{Win probability.}
Let $c$ be the engine evaluation of a position in centipawns from the mover's
side, limited to the range from $-1000$ to $1000$, with a forced mate counted
as the corresponding limit. Lichess maps the evaluation to a win probability
in percent~\cite{lichessaccuracy},
\[
W(c) = 50 + 50\left(\frac{2}{1 + e^{-0.00368208\,c}} - 1\right).
\]
The win probability of an example is computed from the evaluation of the
position before the mover's move. The evaluations are those that Lichess
attaches to a game when a user requests computer analysis.\iflook{} Because
$c$ is limited, the win probability lies between
\V{describe.win-before.min}\% and \V{describe.win-before.max}\%.\fi{} The
curve is symmetric, $W(-c) = 100 - W(c)$, so the win probabilities of the two
players sum to 100\% and no probability is assigned to a draw. We therefore
read $W(c)/100$ as an expected score when we use it as a prediction. The
coaching application uses the same curve, and a parity check in the
repository verifies that the two implementations agree.

\subsection{Dataset}

The source is the Lichess database of rated games~\cite{lichessdb} for
\V{data.month}. A game is included when it is rated, both players have a
rating and neither is a bot, the time control is rapid by Lichess's
definition (an estimated duration of 8 to 25 minutes, where the estimate is
the initial time plus 40 times the increment), the game ended on the board
or on time, and its record carries engine evaluations and clock times.
Lichess computes evaluations for a game only when a user requests an
analysis, so the analysed games are a self-selected subset of all games.

The monthly file is read from its beginning until \V{recipe.games} games have
contributed positions, and \V{manifest.games-read} games were read to find
them. The dataset therefore consists of the first qualifying games of the
month in file order. It is not a random sample of the month.

From each game we take up to \V{recipe.per-game} positions, selected among
the eligible plies by a hash of the game identifier and the ply. A ply is
eligible from ply \V{recipe.min-ply} on, that is after each side's first five
moves, the opening cut used by Maia~\cite{maia}, provided that the
evaluations before and after the move and the clock times around the move are
all available. The dataset is shared with the blunder-risk and move-time
tasks of this repository, which is why the evaluation after the move is
required although no model of this paper uses it. The resulting dataset has
\V{data.rows} positions (Table~\ref{tab:data}).

The split is determined by a hash of the game identifier, so all positions
of a game fall in the same set, and the test set contains positions from
\V{data.games.test} games. All positions of a game share one result, so the
number of games, and not the number of positions, determines how precisely a
score can be estimated. The same grouping is used for the bootstrap
(Section~\ref{sec:protocol}). The dataset is not stored in the repository.
Its recipe is, and a build of the dataset is identified by a fingerprint, a
hash of its contents. All models in this paper were trained and evaluated on
build \V{data.fingerprint}\ifmore; the models of
Section~\ref{sec:training-size} that use more training games were trained on
a larger build and evaluated on the validation and test sets of this
one\fi. The evaluation code refuses to rank models evaluated on different
builds.

\begin{table}[htbp]\centering\small
\block{data}
\caption{Dataset statistics by split. The mean target is the mean result for
the mover.}
\label{tab:data}
\end{table}

\iflook
\paragraph{Distribution of the target.}
The mean result is \V{label.mean}. It is close to one half because the
positions of both players are included, and one player's win is the other's
loss. Draws are rare (Figure~\ref{fig:label}), so the standard deviation of
the result, \V{label.sd}, is close to the 0.5 that it would be without draws.
The standard deviation is the RMSE of the best constant prediction and sets
the scale for every RMSE in this paper. Table~\ref{tab:describe} summarises
the main variables. Half of the positions have a win probability between
\V{describe.win-before.q25}\% and \V{describe.win-before.q75}\%. Its mean is
\V{describe.win-before.mean}\% and not 50\%, which
Section~\ref{sec:engine-baseline} explains. Time pressure is uncommon: the
median time remaining on the mover's clock is \V{describe.clock.q50} seconds,
and it is below \V{describe.clock.q5} seconds in one position in twenty and
below \V{describe.clock.q1} seconds in one in a hundred.

\begin{figure}[htbp]\centering
\block{label}
\caption{Distribution of the target over all positions of the dataset: a
loss, a draw or a win for the mover.}
\label{fig:label}
\end{figure}

\begin{table}[htbp]\centering\small
\block{describe}
\caption{Summary statistics of selected variables over the whole dataset:
the mean, five quantiles, the minimum and the maximum. Clock times are in
seconds and win probabilities in percent. The win probability lost by the
move and the length of the game are known only afterwards and are not
inputs.}
\label{tab:describe}
\end{table}
\fi

\section{Methods}
\label{sec:methods}

\subsection{Input representation}
\label{sec:input}

% Input representation: identical for every task, so each model paper inputs
% this file under its own heading and then adds what is specific to its task.
% The prose follows docs/writing.md.
All tasks in this project share one input representation, so two tasks differ
in their target and in the information a model is given, and not in how a
position is encoded.

\paragraph{Square tokens.}
A position is encoded as 64 tokens, one per square. Token 0 denotes an empty
square, tokens 1 to 6 the pawn, knight, bishop, rook, queen and king of the
side to move, which we call the \emph{mover}, and tokens 7 to 12 the
opponent's pieces. When Black is to move, the board is flipped vertically and
the colours are swapped, so that the mover always plays up the board. A
position with White to move and the corresponding colour-reversed position
with Black to move therefore have the same encoding, and a model does not
have to learn the two cases separately.

\paragraph{Auxiliary input planes.}
Information about individual squares is supplied as binary planes over the 64
squares, which a network receives together with the tokens. This task uses
\V{task.marks.count} planes (Table~\ref{tab:marks}). The first marks the en
passant target square, and the others mark the origin and destination squares
of moves. A plane identifies the squares of a move and carries no other
information about it, such as whether the move is a capture or a check.

\begin{table}[htbp]\centering\small
\block{marks}
\caption{Auxiliary input planes. Squares are given from the mover's point of
view, after the board has been flipped.}
\label{tab:marks}
\end{table}

\ifsequence
\paragraph{Sequences of positions.}
An example of this task is not a single position. It is the sequence of one
player's moves in a game, and each element of the sequence is a position
encoded as above: the board before the move, with the opponent's preceding
move and the move played marked in the auxiliary planes. Scalar inputs
therefore come at two levels, one for each move and one for the game.

\paragraph{Per-move features.}
Each position of the sequence has \V{task.steps.count} scalar features of its
own (Table~\ref{tab:steps}): the four castling rights and the task's
information about that move. A network receives them together with the
position they belong to.

\begin{table}[htbp]\centering\small
\block{steps}
\caption{Per-move features: the scalar inputs attached to each position of a
sequence.}
\label{tab:steps}
\end{table}

\paragraph{Scalar context features.}
Information about the game as a whole is given as \V{task.context.count}
scalar context features (Table~\ref{tab:reads}). Every model receives them,
including the gradient-boosted trees, which never see a board. For a network,
both kinds of scalar feature are standardised with the mean and standard
deviation of the training set.

\begin{table}[htbp]\centering\small
\block{reads}
\caption{Scalar context features: the game-level inputs given to every
model.}
\label{tab:reads}
\end{table}

\paragraph{Hand-crafted features.}
Gradient-boosted trees cannot take a sequence of boards as input, so we also
compute \V{task.hand.count} aggregate features over the player's moves
(Table~\ref{tab:hand} in the appendix): how much win probability the moves
lost on average, at worst and how often beyond fixed thresholds, how the clock
was used, and how many moves were available. These features have two roles.
They are the inputs of the GBDT baseline, and so define the level that a
network reading every position has to exceed. They can also be given to a
network as additional context features, which tests whether the network
already extracts the same information from the sequence.

\else
\paragraph{Scalar context features.}
Information that does not belong to a square is given as scalar context
features: the four castling rights and the scalar inputs specific to the
task. This task uses \V{task.context.count} (Table~\ref{tab:reads}). Every
model receives them, including the gradient-boosted trees, which never see
the board. For a network they are standardised with the mean and standard
deviation of the training set.

\begin{table}[htbp]\centering\small
\block{reads}
\caption{Scalar context features: the inputs given to every model in addition
to the board.}
\label{tab:reads}
\end{table}

\paragraph{Hand-crafted features.}
Gradient-boosted trees cannot take 64 tokens as input, so we also compute
\V{task.hand.count} hand-crafted features of the position
(Table~\ref{tab:hand} in the appendix): material, mobility, available checks
and captures, hanging pieces, passed pawns and attacks near each king, and,
where the task gives a line of moves, features of that line. These features
have two roles. They are the inputs of the GBDT baseline, and so define the
level that a network reading the board has to exceed. They can also be given
to a network as additional context features, which tests whether the network
already extracts the same information from the board.

\fi

\iflook
\begin{figure}[htbp]\centering
\block{encoding}
\ifsequence
\caption{One position of one example, shown twice: the position before the
player's most costly move of the game. Left: as played, with arrows for the
opponent's preceding move (grey) and the player's move (blue). Right: as
encoded, oriented so that the mover plays up the board, with the squares set
in an auxiliary plane shaded grey for the opponent's move and blue for the
player's move. Every position of the sequence is encoded in this way.}
\else
\caption{One example from the dataset, shown twice. Left: as played, with
arrows for the opponent's preceding move (grey) and the move the example
concerns (blue). Right: as encoded, oriented so that the mover plays up the
board, with the squares set in an auxiliary plane shaded grey for the
opponent's move and, where the task marks a line of moves, blue for the
mover's moves and orange for the replies.}
\fi
\label{fig:encoding}
\end{figure}
\fi


\paragraph{Inputs of this task.}
The auxiliary planes mark the en passant square and the origin and
destination squares of the opponent's preceding move. Besides the castling
rights, the scalar context features are the two ratings, the increment, the
ply, three durations, each encoded as the logarithm of one plus the number of
seconds (the time remaining on the mover's clock, the time remaining on the
opponent's clock and the initial time of the time control), and the mover's
win probability before the move. The win probability is the engine's
assessment of the position under best play. It is the one input that is
itself an estimate of the target, and this paper asks what the other inputs
add to it. No input describes the move that is played or anything that
happens after it.\iflook{} In Figure~\ref{fig:encoding} the move played is
drawn on the left board for the reader, and it is not marked on the encoded
board.\fi

\subsection{Models and training}
\label{sec:models}

% Models and training: the same sequence of models for every task, so each
% model paper inputs this file under its own heading.
% The prose follows docs/writing.md.
Every task is given the same sequence of models, so that each step of the
sequence measures the contribution of one component. A model is fitted on the
training set, selected on the validation set and scored once on the test set.
Table~\ref{tab:runs} lists the models of this task.

\paragraph{Constant baseline.}
\run{constant} predicts the same value for every example: the mean of the
training targets (for a binary target, the base rate), or their median when
the primary metric is an absolute error. A useful model has to improve on it.

\paragraph{Gradient-boosted decision trees.}
\ifcontextgbdt
\run{gbm-context} is a gradient-boosted decision tree model (GBDT), trained
with LightGBM~\cite{lightgbm} on the scalar context features alone. It never
sees the board, and so measures how much of the target the context already
determines. \run{gbm-hand} adds the hand-crafted features.
\else
\run{gbm-hand} is a gradient-boosted decision tree model (GBDT), trained with
LightGBM~\cite{lightgbm} on the scalar context features and the hand-crafted
features.
\fi
We refer to \run{gbm-hand} as \emph{the GBDT baseline}: every other model is
compared with it on the same test examples. GBDT is trained for at most
\V{method.gbm.rounds} boosting rounds with up to \V{method.gbm.leaves} leaves
per tree and a learning rate of \V{method.gbm.rate}. Each tree is fitted on a
random \V{method.gbm.rows.percent}\% of the examples and
\V{method.gbm.features.percent}\% of the features, and boosting stops when the
validation score has not improved for \V{method.gbm.patience} rounds.

\paragraph{Neural networks.}
A network consists of an encoder, which maps the square tokens, the auxiliary
planes and the context features to a vector, and a prediction head, a
two-layer perceptron that maps this vector to the output. We use two encoder
architectures, so that the results show how much the architecture matters for
the task.

The \emph{transformer} encoder~\cite{attention} treats each square as a token.
The embedding of a token is the sum of a learned embedding of its piece, a
learned embedding of its square and a linear projection of the square's
auxiliary planes. The context features are projected by a small perceptron
and enter as one additional token. Every token attends to every other token
in every layer, so distant squares interact from the first layer on. To give
attention the geometry of the board, each attention head adds a learned bias
that depends on the file and rank offset between two
squares~\cite{relative}: files, ranks, diagonals and knight moves all
correspond to fixed offsets. The layers use pre-layer normalisation, and the
encoder outputs the context token concatenated with the mean of the 64 square
tokens.

The \emph{residual convolutional network} (ResNet) is a stack of residual
blocks of $3\times3$ convolutions over the $8\times8$ board, the architecture
that AlphaZero introduced for chess~\cite{alphazero}. Its input is one binary
plane per token value together with the auxiliary planes. A linear projection
of the context features is added at every square after the first convolution,
and the encoder outputs the mean and the maximum over the squares of the
final feature map. Unlike attention, a convolution relates only neighbouring
squares, so the receptive field grows with depth.

\ifsequence
\paragraph{Sequence model.}
An example of this task is a sequence of positions, so a network has two
stages below its prediction head. One of the two encoders above is applied to
every position of the sequence, with the position's per-move features in
place of the context features, and produces one vector per position. A
two-layer transformer then operates on these vectors in the order in which
the moves were played: each vector is a token, a learned positional embedding
is added to it, and the game-level context features enter as one additional
token. The prediction head receives that token concatenated with the mean of
the position tokens. Sequences in a batch are padded to the length of the
longest, and padded elements are masked: they are not encoded and no token
attends to them, so a prediction does not depend on the other examples in
the batch.

Model sizes are hyperparameters. \run{conv-small} is small enough to train
on a CPU. An example consists of tens of positions, so
\run{transformer-small} and \run{transformer-medium} are trained on a GPU.
\run{transformer-medium-hand} is \run{transformer-medium} with the
hand-crafted features added to its context features.
Table~\ref{tab:settings} in the appendix lists the hyperparameters.
\else
Model sizes are hyperparameters. \run{conv-small} and \run{transformer-small}
are small enough to train on a CPU. \run{transformer-medium} is trained on a
GPU, and \run{transformer-medium-hand} is the same network with the
hand-crafted features added to its context features.
Table~\ref{tab:settings} in the appendix lists the hyperparameters.
\fi

\paragraph{Training.}
All networks are trained in the same way: with AdamW~\cite{adamw}, a learning
rate that increases linearly over the first \V{networks.warmup.percent}\% of
the optimisation steps and then follows a cosine schedule to zero, gradient
clipping, and mixed precision on a GPU. The loss matches the primary metric:
mean squared error for a target scored by RMSE, a Huber loss for a target
scored by MAE, and binary cross-entropy for a probability. A regression
target is standardised with the mean and standard deviation of the training
set. After every epoch the network is evaluated on the validation set. We
keep the weights of the best epoch, and stop training when
\V{networks.patience} consecutive epochs have not improved on it.
\ifmore
A model trained on a larger training set is given proportionally fewer
epochs, and a correspondingly shorter patience
(Section~\ref{sec:ablations}).
\fi

\begin{table}[htbp]\centering\footnotesize
\block{runs}
\caption{The models, best first. \emph{Selected} is the number of trees or
the training epoch chosen on the validation set. \emph{Time} is the total
wall-clock time for training and evaluation.}
\label{tab:runs}
\end{table}


\paragraph{Baselines without fitted parameters.}
Two further baselines are computed from the inputs without any training.
\run{engine-curve} is the engine-derived win probability used directly as
the prediction, $\hat{y} = W(c)/100$. It is the value that analysis software
displays, and the first model that a fitted model has to improve on.
\run{elo-gap} uses only the mover's rating $r_m$ and the opponent's rating
$r_o$. Its prediction is the expected score of the Elo formula,
\[
\hat{y} = \frac{1}{1 + 10^{-(r_m - r_o)/400}},
\]
and it shows how much of the result the two ratings predict without any
knowledge of the position.

\section{Evaluation Protocol}
\label{sec:protocol}

\paragraph{Metrics.}
The primary metric, fixed before the first full training run, is the root
mean squared error of the predicted score on the test set, where lower is
better. We chose it because the squared error is a proper scoring rule for an
expectation: its expected value is smallest for the expected score and for no
other prediction. A model that predicts 0.7 for positions from which movers
score 0.7 on average cannot be improved on those positions. We report three
secondary metrics. The coefficient of determination R\textsuperscript{2} is
the share of the variance of the result that a model explains, and
Spearman's rank correlation between the predicted score and the result
measures only how well a model orders the positions. The mean absolute error
(MAE) is reported for completeness and should not be used to rank models of
an expected score. The expected absolute error is smallest for the median of
the result, which is 1 whenever the mover wins more than half of such games
and 0 whenever the mover loses more than half, and not for the expected
score. The MAE therefore favours predictions that are more extreme than the
expected score. We also report results separately for \V{bands.count} bands
of the mover's rating.

% Evaluation protocol: the same for every task. Each model paper names its
% own metrics under its section heading and inputs the rest from here.
% The prose follows docs/writing.md.
\paragraph{Confidence intervals.}
Examples from the same game are not independent, so we resample whole games.
The test games are drawn with replacement \V{method.resamples} times, the
metric is computed each time on the examples of the games drawn, and the
2.5th and 97.5th percentiles of these values form the 95\% confidence
interval (CI) that we report with every score.

\paragraph{Paired comparisons.}
To compare two models we compute both scores on the same resampled games and
take their difference. We call this difference the \emph{gain} of the first
model over the second, and sign it so that a positive gain means that the
first model is better. A gain is \emph{statistically significant} when its
95\% CI excludes zero, and is then set in bold in the tables. Unless stated
otherwise, the second model is the GBDT baseline. Two models tend to do well
and badly on the same games, so a paired gain can be significant even when
the two models' separate intervals overlap. These intervals reflect the
sampling of test games only. Each network was trained with a single random
seed, and the variation between training runs is not included.

\paragraph{Model selection.}
\ifchance
The number of trees, the training epoch whose weights are kept and the
calibration parameters are chosen on the validation set.
\else
The number of trees and the training epoch whose weights are kept are chosen
on the validation set.
\fi
The test set is used once per model, for the scores reported here.

\ifchance
\paragraph{Calibration.}
Before a model is scored, its predicted probabilities are rescaled by Platt
scaling~\cite{platt}: a logistic function of the log-odds whose slope and
intercept are fitted on the validation set. The transformation is monotone,
so it leaves a model's ranking of the examples unchanged and affects only its
log loss and its calibration. What remains of a model's miscalibration after
this step cannot be removed by a linear correction of its log-odds.
\fi

\ifpairs
\paragraph{Acceptance criterion.}
Section~\ref{sec:ablations} evaluates modifications of the models. A
modification is \emph{accepted} only if the modified model is better than the
unmodified one on the validation set and its paired gain on the test set is
statistically significant. A modification that is not accepted is still
reported, as a measurement of its effect.
\fi

\paragraph{Tables and figures.}
A model has the same colour and marker in every figure, and the GBDT baseline
is always drawn in orange. The values plotted in a figure are also given in a
table, in the text or in the appendix. All numbers, tables and figures are
generated from the committed result files by the command that builds the
paper.


\section{Results}
\label{sec:results}

\subsection{Main results}

Table~\ref{tab:leaderboard} and Figure~\ref{fig:scores} give the main
results. The constant baseline predicts the mean result of the training set,
\V{floor.constant}, for every position and has an RMSE of \V{floor.score}.
The engine-derived win probability, used directly as a prediction, reaches
\V{run.engine-curve.score}. Without any fitted parameter it achieves
\V{run.engine-curve.lift.share}\% of the improvement of the GBDT baseline
over the constant baseline, so the evaluation of the position carries most of
what the models predict. GBDT on the context features alone, which are the
win probability, the ratings, the clocks, the time control, the ply and the
castling rights, reaches \V{run.gbm-context.score} and achieves
\V{context.lift.share}\% of that improvement. The GBDT baseline, with the
hand-crafted features of the position, reaches \V{reference.score}. The Elo
expected score computed from the two ratings has an RMSE of
\V{run.elo-gap.score} and is the only model that is less accurate than the
constant baseline. Sections~\ref{sec:engine-baseline}
and~\ref{sec:rating-baseline} examine the two baselines without fitted
parameters.

\begin{table}[htbp]\centering\small
\block{leaderboard}
\caption{Main results on the \V{data.rows.test} test positions, from
\V{data.games.test} games, best first. The gain is the paired difference in
RMSE from \V{task.reference} on the same resampled games. It is positive when
the model is better, and bold when statistically significant.}
\label{tab:leaderboard}
\end{table}

\begin{figure}[htbp]\centering
\block{scores}
\caption{Test RMSE of every model with its 95\% CI. The vertical line marks
the GBDT baseline.}
\label{fig:scores}
\end{figure}

The networks of the fixed sequence do not improve on the GBDT baseline. None
of the \V{ladder.count} differs significantly from it
(Figure~\ref{fig:gains}). The gains are \V{run.conv-small.gain} for the small
ResNet (95\% CI \V{run.conv-small.gain.low} to \V{run.conv-small.gain.high}),
\V{run.transformer-small.gain} for the small transformer
(\V{run.transformer-small.gain.low} to \V{run.transformer-small.gain.high}),
\V{run.transformer-medium.gain} for the medium transformer
(\V{run.transformer-medium.gain.low} to \V{run.transformer-medium.gain.high})
and \V{run.transformer-medium-hand.gain} for the medium transformer with the
hand-crafted features (\V{run.transformer-medium-hand.gain.low} to
\V{run.transformer-medium-hand.gain.high}). Each of these intervals is wider
than the largest of the four gains in absolute value, so the test set does
not show whether any of these networks is better or worse than GBDT.
\ifpairs
The hand-crafted features do not improve the medium transformer: compared
with \run{transformer-medium}, \run{transformer-medium-hand} is
\V{pair.transformer-medium-hand.transformer-medium.verdict} on the test set
(gain \V{pair.transformer-medium-hand.transformer-medium.gain}; 95\% CI
\V{pair.transformer-medium-hand.transformer-medium.gain.low} to
\V{pair.transformer-medium-hand.transformer-medium.gain.high}).
\fi

\begin{figure}[htbp]\centering
\block{gains}
\caption{Paired gain in RMSE of each network over the GBDT baseline, with its
95\% CI. A model whose interval lies entirely to the right of the line is
significantly better than the baseline, and one whose interval lies entirely
to the left is significantly worse.}
\label{fig:gains}
\end{figure}

The remaining models in Table~\ref{tab:leaderboard} are modifications that
passed the acceptance criterion of Section~\ref{sec:protocol}\ifmore: GBDT
and the medium transformer trained on larger training sets
(Section~\ref{sec:training-size})\fi. The best model is \V{best.name}, the
GBDT baseline trained on \V{more.gbm-hand.most} the training games, with an
RMSE of \V{best.score} (95\% CI \V{best.score.low} to \V{best.score.high}).
Its gain over the GBDT baseline is \V{best.gain} (95\% CI \V{best.gain.low}
to \V{best.gain.high}), which is \V{best.gain.percent}\% of the baseline's
RMSE and \V{best.gain.of-lift}\% of the baseline's improvement over the
constant baseline. The best network, \V{network.name}, has an RMSE of
\V{network.score} and does not differ significantly from the GBDT baseline
(gain \V{network.gain}; 95\% CI \V{network.gain.low} to
\V{network.gain.high}). Compared with the GBDT baseline,
\V{networks.ahead.of} on the leaderboard are significantly better and
\V{networks.behind.of} are significantly worse.

\subsection{The engine-derived win probability as a predictor}
\label{sec:engine-baseline}

\run{engine-curve} has an RMSE of \V{run.engine-curve.score} (95\% CI
\V{run.engine-curve.score.low} to \V{run.engine-curve.score.high}).
\ifpairs
\run{gbm-context} receives the same win probability together with the
ratings, the clocks, the time control and the ply, and does not see the
board. Compared with \run{engine-curve} it is
\V{pair.gbm-context.engine-curve.verdict} on the test set (gain
\V{pair.gbm-context.engine-curve.gain}; 95\% CI
\V{pair.gbm-context.engine-curve.gain.low} to
\V{pair.gbm-context.engine-curve.gain.high}), and the gain is
\V{pair.gbm-context.engine-curve.gain.percent}\% of the RMSE of the win
probability. The medium transformer with hand-crafted features is
\V{pair.transformer-medium-hand.engine-curve.verdict} as well (gain
\V{pair.transformer-medium-hand.engine-curve.gain}; 95\% CI
\V{pair.transformer-medium-hand.engine-curve.gain.low} to
\V{pair.transformer-medium-hand.engine-curve.gain.high};
Table~\ref{tab:pairs}).
\fi
Relative to the GBDT baseline, the gain of \run{engine-curve} is
\V{run.engine-curve.gain} (95\% CI \V{run.engine-curve.gain.low} to
\V{run.engine-curve.gain.high}).

The gains of the fitted models are small relative to the error from which
they are taken, and they cannot be large. From a single position most of the
result is still undecided, and no model in this paper explains more than
\V{best.r2} of the variance of the result (Table~\ref{tab:metrics}). The
fitted models do not predict individual games much better than the win
probability does. Their advantage is calibration: their predictions agree
with the mean result over the range of predictions and across ratings and
clocks, and the win probability does not.

\iflook
\paragraph{Calibration.}
In Figure~\ref{fig:reliability} the test positions are divided into
\V{reference.reliability.groups} bins of equal size by a model's prediction,
and the mean result of each bin is plotted against the mean prediction. For
the GBDT baseline the two differ by at most
\V{reference.reliability.widest-gap} in any bin. The win probability deviates
at both ends. In its lowest bin it gives the mover
\V{run.engine-curve.reliability.lowest.said} on average and movers score
\V{run.engine-curve.reliability.lowest.happened}. In its highest bin it gives
\V{run.engine-curve.reliability.highest.said} and movers score
\V{run.engine-curve.reliability.highest.happened}. The largest difference,
\V{run.engine-curve.reliability.widest-gap}, occurs in a bin in which it
gives \V{run.engine-curve.reliability.widest-gap.said} and movers score
\V{run.engine-curve.reliability.widest-gap.happened}. The win probability is
thus overconfident in both directions: movers who stand badly score more than
it indicates, and movers who stand well score less.
Section~\ref{sec:error} shows how these deviations depend on the rating and
on the clock.

\begin{figure}[htbp]\centering
\block{reliability-engine-curve}
\caption{Mean result against mean prediction for the engine-derived win
probability and the GBDT baseline. The test positions are divided into bins
of equal size by a model's prediction, and each point is one bin. Points on
the diagonal are unbiased.}
\label{fig:reliability}
\end{figure}

\paragraph{Offset of an evaluation taken before the move.}
The deviations are not symmetric. Over the whole test set the mean of the win
probability is \V{overall.engine-curve.said} and the mean result is
\V{overall.happened}, so the mean prediction minus the mean result is
\V{overall.engine-curve.off}. For the GBDT baseline this difference is
\V{overall.gbm-hand.off}. The offset follows from the moment at which the
evaluation is taken. Let $b_k$ be the mover's win probability before move $k$
of a game, in percent, and $l_k$ the win probability that the move loses.
After the move the mover's win probability is $b_k - l_k$, and the opponent's
win probability before the next move is $b_{k+1} = 100 - (b_k - l_k)$. Hence
$l_k = b_k + b_{k+1} - 100$, and averaged over the moves of a game the mean
win probability before a move is 50\% plus half the mean loss per move. An
engine evaluation assumes best play by both sides, so the moves that players
actually make lose win probability on average. In the dataset the mean loss
is \V{describe.loss.mean} percentage points and the mean win probability
before the move is \V{describe.win-before.mean}\%, in agreement with this
relation (Table~\ref{tab:describe}).

The win probability before a move therefore overstates the mover's expected
score on average, by half of what a move loses. When the evaluation is taken,
the opponent has just made a move that was not perfect on average, and the
mover has not yet made one. An application that displays the win probability
before a move shows a value that is too favourable to the player to move by
this amount on average.
\fi

\subsection{The two ratings as a predictor}
\label{sec:rating-baseline}

The Elo expected score computed from the two ratings, \run{elo-gap}, has an
RMSE of \V{run.elo-gap.score} (95\% CI \V{run.elo-gap.score.low} to
\V{run.elo-gap.score.high}), which is higher than the \V{floor.score} of the
constant baseline. Its R\textsuperscript{2} is negative (\V{run.elo-gap.r2};
95\% CI \V{run.elo-gap.r2.low} to \V{run.elo-gap.r2.high}), which means that
its squared error is larger than that of predicting the mean result of the
test set for every position. Its rank correlation with the result is
\V{run.elo-gap.spearman} (Table~\ref{tab:metrics}).
\ifpairs
\run{gbm-context}, whose inputs include the two ratings, is
\V{pair.gbm-context.elo-gap.verdict} than \run{elo-gap} on the test set (gain
\V{pair.gbm-context.elo-gap.gain}; 95\% CI
\V{pair.gbm-context.elo-gap.gain.low} to
\V{pair.gbm-context.elo-gap.gain.high}).
\fi

\iflook
Figure~\ref{fig:formula} shows the reason. Across the bins the mean result
rises with the formula's prediction, but far less steeply than the diagonal. In the lowest
of the \V{reference.reliability.groups} bins the formula gives the mover
\V{run.elo-gap.reliability.lowest.said} on average and movers score
\V{run.elo-gap.reliability.lowest.happened}. In the highest bin it gives
\V{run.elo-gap.reliability.highest.said} and movers score
\V{run.elo-gap.reliability.highest.happened}. The higher-rated player thus
scores more than one half, but much less than the formula predicts, and the
squared error of the formula's confident predictions outweighs what its
ordering of the positions gains.

\begin{figure}[htbp]\centering
\block{reliability-elo-gap}
\caption{Mean result against mean prediction for the Elo expected score and
the GBDT baseline, in bins of equal size by a model's prediction. Points on
the diagonal are unbiased.}
\label{fig:formula}
\end{figure}
\fi

The formula is not equally inaccurate at all ratings. Its RMSE is higher than
that of the constant baseline in each of the four bands of the mover's rating
below 2000, and lower in the band \V{band.5.name} (\V{run.elo-gap.band.5}
against \V{band.5.floor}; Table~\ref{tab:bands}).

We have not identified the cause of the miscalibration. That a position is
more informative than the ratings does not explain it. The formula describes the result of a game,
and every position of a game has the result of that game, so the calibration
of the formula on these positions is its calibration on the games themselves,
up to the weighting of a game by its number of positions. Two explanations
are possible, and the dataset does not contain the variables needed to test
them. First, a large rating difference may often involve a rating that is not
yet reliable, such as that of a new account.\iflook{} Most games in the
dataset are between players of similar rating: only in the lowest and the
highest bin of Figure~\ref{fig:formula} is the formula's prediction far from
one half.\fi{} Glicko-2 records the uncertainty of a rating separately, as
the rating deviation, and the Elo formula applied to the rating alone ignores
it. Second, the dataset contains only games for which a user requested
computer analysis, and games with an unexpected result may be analysed more
often.\iflook{} The largest errors of the best model are consistent with the
first explanation without establishing it. They include a game in which a
player rated exactly 1500, the rating that Lichess assigns to a new account,
defeated an opponent rated more than 900 points higher
(Section~\ref{sec:positions}).\fi

This result does not mean that the ratings are uninformative. The
opponent's rating and the mover's rating are the second and third most
important features of the GBDT baseline
(Section~\ref{sec:importance})\iflook, and the relation between the win
probability and the result depends strongly on the mover's rating
(Section~\ref{sec:error})\fi.

\subsection{Secondary metrics}

Table~\ref{tab:metrics} gives all metrics for every model. The
R\textsuperscript{2} of the best model is \V{best.r2}, compared with
\V{reference.r2} for the GBDT baseline, \V{run.gbm-context.r2} for GBDT on
the context features alone and \V{run.engine-curve.r2} for the engine-derived
win probability. The rank correlations with the result are \V{best.spearman},
\V{reference.spearman}, \V{run.gbm-context.spearman} and
\V{run.engine-curve.spearman}. RMSE, R\textsuperscript{2} and the rank
correlation order these four models in the same way.

\begin{table}[htbp]\centering\small
\block{metrics}
\caption{All metrics for every model on the test set. RMSE and MAE are errors
of the predicted score and lower is better. The models are ranked by RMSE;
the MAE favours extreme predictions (Section~\ref{sec:protocol}). For
R\textsuperscript{2} and Spearman's rank correlation higher is better. The
constant baseline predicts a single value and has no rank correlation.}
\label{tab:metrics}
\end{table}

The MAE orders the models differently. The engine-derived win probability has
the lowest MAE of all models (\V{run.engine-curve.mae}, compared with
\V{reference.mae} for the GBDT baseline), although every fitted model has a
lower RMSE. This is the property of the absolute error described in
Section~\ref{sec:protocol}. The win probability is more extreme than the
expected score, which lowers its absolute error and raises its squared
error.\iflook{} For the same reason its prediction is closer to the result
than that of the GBDT baseline on \V{run.engine-curve.beside.closer}\% of the
test positions (Section~\ref{sec:error}).\fi

\subsection{Results by rating band}

Table~\ref{tab:bands} gives the RMSE by band of the mover's rating. The RMSE
of the constant baseline is between \V{band.4.floor} and \V{band.2.floor} in
the four bands below 2000 and \V{band.5.floor} in the band \V{band.5.name}.
It measures the spread of the results in a band, so scores should be compared
between models within a band and not between bands. Every model has its
lowest RMSE in the highest band.

\begin{table}[htbp]\centering\small
\block{bands-table}
\caption{Test RMSE by band of the mover's rating. The spread of the results
differs between bands, so values are comparable within a column and not
between columns.}
\label{tab:bands}
\end{table}

The best model has a lower RMSE than the GBDT baseline in
\V{bands.best-ahead.of}. The difference is largest in the band
\V{bands.lead.widest.name} (\V{bands.lead.widest}) and smallest in the band
\V{bands.lead.narrowest.name} (\V{bands.lead.narrowest}).\iflook{} With a
paired comparison within each band, the gain is statistically significant in
\V{slice.band.ahead.of} (Table~\ref{tab:slice-band} in the appendix). The
engine-derived win probability has a higher RMSE than the GBDT baseline in
every band. The difference is largest in the band
\V{slice.band.engine-curve.behind.most.group}
(\V{slice.band.engine-curve.behind.most}) and smallest in the band
\V{slice.band.engine-curve.behind.least.group}
(\V{slice.band.engine-curve.behind.least}). These are differences between two
scores on the same positions, without confidence intervals.\fi

\subsection{Training curves}

Figure~\ref{fig:curves} shows the validation RMSE of each network of the
fixed sequence after every epoch. By this measure
\V{ladder.at-last-epoch.of} were still improving when training ended.
\run{conv-small} had its best validation RMSE at epoch
\V{run.conv-small.best-epoch} of \V{run.conv-small.epochs} and
\run{transformer-small} at epoch \V{run.transformer-small.best-epoch} of
\V{run.transformer-small.epochs}. \run{transformer-medium} had it at epoch
\V{run.transformer-medium.best-epoch} of \V{run.transformer-medium.epochs}
and \run{transformer-medium-hand} at epoch
\V{run.transformer-medium-hand.best-epoch}, and early stopping ended the
training of both after epoch \V{run.transformer-medium.epochs.run}.
\ifmore
The medium transformer trained on \V{more.transformer-medium-hand.most} the
games is the exception among the networks of this paper: its best epoch was
its last (Section~\ref{sec:training-size}).
\fi

\begin{figure}[htbp]\centering
\block{curves}
\caption{Validation RMSE of each network of the fixed sequence after every
epoch. The dashed line is the GBDT baseline's validation RMSE.}
\label{fig:curves}
\end{figure}

\subsection{Feature importance of the GBDT baseline}
\label{sec:importance}

The GBDT baseline consists of \V{reference.trees} trees over
\V{reference.features} features. Figure~\ref{fig:leans} shows the ten
features with the largest shares of the total split gain. The win probability
before the move accounts for \V{reference.leans.1.share}\% of the gain. It is
followed by the \V{reference.leans.2.feature}
(\V{reference.leans.2.share}\%), the \V{reference.leans.3.feature}
(\V{reference.leans.3.share}\%), the \V{reference.leans.4.feature}
(\V{reference.leans.4.share}\%), the \V{reference.leans.5.feature}
(\V{reference.leans.5.share}\%) and the \V{reference.leans.6.feature}
(\V{reference.leans.6.share}\%). The model thus relies mainly on the engine's
assessment and adjusts it for the ratings, the material and the clocks. Split
gain shows which features the model used. It does not identify what decides a
game, because correlated features divide the credit between them arbitrarily.

\begin{figure}[htbp]\centering
\block{leans}
\caption{The ten most important features of the GBDT baseline, by share of
the total split gain.}
\label{fig:leans}
\end{figure}

\ifchanges
\section{Ablations}
\label{sec:ablations}

% Ablations: what a modification is, the table of paired comparisons and the
% two figures of validation curves. The same for every task, so a model paper
% inputs this file at the start of its ablation section and then reports its
% own findings. Each part appears once the task has the runs it needs, and the
% list below needs at least one of its items.
% The prose follows docs/writing.md.
\ifpairs
The main results compare complete models. An ablation isolates one factor: it
compares a model with its \emph{base}, a model that differs from it in a
single respect. Table~\ref{tab:pairs} lists every such comparison for this
task, with the gain of the modified model over its base on the validation set
and its paired gain on the test set, computed on the same resampled test
games as in the main results. The acceptance criterion of
Section~\ref{sec:protocol} is applied to these pairs. The modifications are
the following.
\begin{itemize}
\ifchangehand
\item \textbf{Hand-crafted features} (suffix \code{-hand}): the network
receives the hand-crafted features as additional context features.
\fi
\ifchangeepochs
\item \textbf{Longer schedule} (suffix \code{-long}): twice the training
epochs, with the same warm-up fraction and cosine schedule.
\fi
\ifchangestart
\item \textbf{Pretraining} (suffix \code{-pretrained}): the transformer
encoder is initialised from a transformer of the same size that was trained
to predict the engine's win probability for positions from the Lichess
evaluation database~\cite{lichessdb}. The piece and square embeddings and all
transformer layers are copied. The prediction head is initialised randomly,
as are the projections of the auxiliary planes that the pretraining task does
not have and, unless the two tasks have the same context features, the
projection of the context features. Training then follows the standard
schedule.
\fi
\ifchangeattacks
\item \textbf{Attack-aware attention bias} (suffix \code{-attacks}): the
relative position bias encodes where two squares lie and not what stands
between them. With this modification each attention head learns four
additional biases for a pair of squares, applied when the piece on one square
attacks the other, when it is attacked from the other, and for the same two
relations through exactly one intervening piece, as in a pin or a skewer. The
relations are computed from the tokens without search, and the biases are
initialised to zero.
\fi
\ifchangegeometry
\item \textbf{No relative position bias} (suffix \code{-no-geometry}): the
relative position bias is removed. This configuration measures the
contribution of the bias and is not a candidate model.
\fi
\ifchangeconv
\item \textbf{ResNet encoder} (prefix \code{conv-medium}): a ResNet with
squeeze-and-excitation in place of the transformer, trained at the higher
peak learning rate used for the convolutional models.
\fi
\ifchangehybrid
\item \textbf{Hybrid encoder} (prefix \code{hybrid-medium}): a few residual
convolutional blocks followed by transformer layers over the resulting
square representations.
\fi
\ifchangelarger
\item \textbf{Larger transformer} (prefix \code{transformer-large}): more layers and
a larger width, with the training schedule unchanged.
\fi
\ifgames
\item \textbf{Training-set size}: the same model trained on a fraction or a
multiple of the training games (Section~\ref{sec:training-size}).
\fi
\iffeatures
\item \textbf{Feature groups}: a model trained without one group of
hand-crafted features at a time (Section~\ref{sec:feature-groups}).
\fi
\end{itemize}

\longblock{pairs}{Paired comparisons. Each row compares a model with a
baseline, in most rows the model it modifies. \emph{Change} states how the
two differ. The gain is that of the model over the baseline and is positive
when the model is better. On
the validation set it is the difference between the two scores. On the test
set it is the paired gain, with its 95\% CI, and is set in bold when
statistically significant.}{tab:pairs}
\fi

\ifvariants
Figure~\ref{fig:variants} shows the validation curves of each model of the
fixed sequence whose training schedule or initialisation was modified,
together with the modified models. A modified model has the colour of its
base and a broken line: long dashes for a longer schedule and dots for a
pretrained encoder.

\begin{figure}[htbp]\centering
\block{variants}
\caption{Validation curves of a model of the fixed sequence (solid line) and
of the models that modify its training schedule or initialisation (broken
lines), after every epoch. One panel per base model, on a common scale. The
dashed horizontal line is the GBDT baseline's validation score.}
\label{fig:variants}
\end{figure}
\fi

\ifshapes
Figure~\ref{fig:shapes} shows the validation curves of the models that
change the architecture and nothing about the training data, each in its own
panel beside its base.

\begin{figure}[htbp]\centering
\block{shapes}
\caption{Validation curves of a model with a modified architecture (broken
line) and of its base (solid line), after every epoch. One panel per modified
model, on a common scale. The dashed horizontal line is the GBDT baseline's
validation score.}
\label{fig:shapes}
\end{figure}
\fi


\ifchangehand
\subsection{Hand-crafted features}

The hand-crafted features give GBDT a small, significant gain and give the
medium transformer none. Compared with \run{gbm-context}, \run{gbm-hand} is
\V{pair.gbm-hand.gbm-context.valid.verdict} than its base on the validation
set and \V{pair.gbm-hand.gbm-context.verdict} on the test set (gain
\V{pair.gbm-hand.gbm-context.gain} RMSE; 95\% CI
\V{pair.gbm-hand.gbm-context.gain.low} to
\V{pair.gbm-hand.gbm-context.gain.high}), so the modification is
\V{pair.gbm-hand.gbm-context.kept}. Compared with \run{transformer-medium},
\run{transformer-medium-hand} is
\V{pair.transformer-medium-hand.transformer-medium.valid.verdict} than its
base on the validation set and
\V{pair.transformer-medium-hand.transformer-medium.verdict} on the test set
(gain \V{pair.transformer-medium-hand.transformer-medium.gain}; 95\% CI
\V{pair.transformer-medium-hand.transformer-medium.gain.low} to
\V{pair.transformer-medium-hand.transformer-medium.gain.high}), so the
modification is \V{pair.transformer-medium-hand.transformer-medium.kept}.
\iffeatures
Section~\ref{sec:feature-groups} shows which groups of features GBDT uses.
\fi
\fi

\ifchangestart
\subsection{Pretraining}

A pretrained encoder makes both medium transformers worse. Without
hand-crafted features, \run{transformer-medium-pretrained} is
\V{pair.transformer-medium-pretrained.transformer-medium.valid.verdict} than
its base on the validation set and
\V{pair.transformer-medium-pretrained.transformer-medium.verdict} on the
test set (gain \V{pair.transformer-medium-pretrained.transformer-medium.gain}
RMSE; 95\% CI
\V{pair.transformer-medium-pretrained.transformer-medium.gain.low} to
\V{pair.transformer-medium-pretrained.transformer-medium.gain.high}), and
the modification is
\V{pair.transformer-medium-pretrained.transformer-medium.kept}. With
hand-crafted features, \run{transformer-medium-hand-pretrained} is
\V{pair.transformer-medium-hand-pretrained.transformer-medium-hand.valid.verdict}
than its base on the validation set and
\V{pair.transformer-medium-hand-pretrained.transformer-medium-hand.verdict}
on the test set (gain
\V{pair.transformer-medium-hand-pretrained.transformer-medium-hand.gain};
95\% CI
\V{pair.transformer-medium-hand-pretrained.transformer-medium-hand.gain.low}
to
\V{pair.transformer-medium-hand-pretrained.transformer-medium-hand.gain.high}),
and the modification is
\V{pair.transformer-medium-hand-pretrained.transformer-medium-hand.kept}.

Both pretrained models had their best validation RMSE earlier than their
bases: \run{transformer-medium-pretrained} at epoch
\V{run.transformer-medium-pretrained.best-epoch} and
\run{transformer-medium-hand-pretrained} at epoch
\V{run.transformer-medium-hand-pretrained.best-epoch} of
\V{run.transformer-medium-pretrained.epochs}, compared with epoch
\V{run.transformer-medium.best-epoch} for both of their bases
(Figure~\ref{fig:variants}). Each was stopped \V{networks.patience} epochs
after its best epoch by early stopping.

Pretraining with the same encoder is accepted in the puzzle-difficulty,
blunder-risk, move-time and tactic-found tasks of this repository. In two of
them, blunder risk and tactics found, the models also receive an
engine-derived win probability as an input, so receiving the evaluation does
not by itself explain the result. What distinguishes the present task is that
the evaluation is close to the target itself: used directly as a prediction,
it achieves \V{run.engine-curve.lift.share}\% of the GBDT baseline's
improvement over the constant baseline. We hypothesise that an encoder
trained to reproduce the engine's evaluation adds little to a model that
already receives this evaluation and predicts a closely related quantity, and
that it then fits the training games sooner than an encoder trained from
random initialisation, as the earlier best epochs suggest. We have not tested
this hypothesis.
\fi

% Training-set size: the same ablation for every task, so a model paper
% inputs this file inside its ablation section and then reports its own
% findings. It appears once the task has the runs.
% The prose follows docs/writing.md.
\ifgames
\subsection{Training-set size}
\label{sec:training-size}

To estimate how much a model's accuracy depends on the amount of training
data, we retrain it on a half and on a quarter of the training games and
evaluate it on the full test set (Table~\ref{tab:games},
Figure~\ref{fig:games}). The subsets are selected by a hash of the game
identifier and are nested, so the quarter is contained in the half. A network
is given twice the epochs on half the games and four times the epochs on a
quarter, so the number of optimisation steps is constant and only the amount
of data varies.
\ifmore
We also train on more data. A larger sample of the same source is split in
the same way as the standard dataset, so the validation and test sets are
unchanged and no test game appears in the larger training set. A network
trained on a larger set is given proportionally fewer epochs, so that it
again takes the same number of optimisation steps as its base. Unlike a model
trained on a subset, a model trained on more data is a candidate model, and
the acceptance criterion applies to it.
\fi

\begin{table}[htbp]\centering\footnotesize
\block{study-games}
\caption{Training-set size. Each model is trained on a fraction of the
training games\ifmore, or on a multiple of them from the larger sample,\fi{}
and evaluated on the full test set. The gain is paired with the same model
trained on the standard training set. A negative gain means that the model is
worse, and bold marks statistical significance.}
\label{tab:games}
\end{table}

\begin{figure}[htbp]\centering
\block{study-games-figure}
\caption{Test score as a function of the number of training games. Each step
to the right doubles the training set.}
\label{fig:games}
\end{figure}
\fi


\ifgames
\paragraph{Smaller training sets.}
GBDT loses little on smaller training sets, and the transformer loses more.
Relative to the same model trained on the full training set, \run{gbm-hand}
has a paired gain of \V{pair.gbm-hand-half.gbm-hand.gain} RMSE on half of the
training games (\V{pair.gbm-hand-half.gbm-hand.verdict}) and of
\V{pair.gbm-hand-quarter.gbm-hand.gain} on a quarter
(\V{pair.gbm-hand-quarter.gbm-hand.verdict}). For
\run{transformer-medium-hand} the gains are
\V{pair.transformer-medium-hand-half.transformer-medium-hand.gain} and
\V{pair.transformer-medium-hand-quarter.transformer-medium-hand.gain}, and it
is significantly worse in \V{study.games.transformer-medium-hand.worse.of}.
On both subsets the transformer has the higher RMSE of the two models
(Table~\ref{tab:games}).
\fi

\ifmore
\paragraph{Larger training sets.}
The larger sample contains up to \V{more.gbm-hand.most} the training games,
with the validation and test sets unchanged. Relative to the same model on
the standard training set, \run{transformer-medium-hand} has a paired gain of
\V{pair.transformer-medium-hand-double.transformer-medium-hand.gain} RMSE on
twice the training games
(\V{pair.transformer-medium-hand-double.transformer-medium-hand.verdict};
\V{pair.transformer-medium-hand-double.transformer-medium-hand.kept}) and of
\V{pair.transformer-medium-hand-quadruple.transformer-medium-hand.gain} on
four times
(\V{pair.transformer-medium-hand-quadruple.transformer-medium-hand.verdict};
\V{pair.transformer-medium-hand-quadruple.transformer-medium-hand.kept}).
For GBDT the gains are \V{pair.gbm-hand-double.gbm-hand.gain}
(\V{pair.gbm-hand-double.gbm-hand.verdict};
\V{pair.gbm-hand-double.gbm-hand.kept}) and
\V{pair.gbm-hand-quadruple.gbm-hand.gain}
(\V{pair.gbm-hand-quadruple.gbm-hand.verdict};
\V{pair.gbm-hand-quadruple.gbm-hand.kept}).

The RMSE of both models decreases with every doubling of the training set,
from a quarter of the games to four times (Figure~\ref{fig:games}). More data
helps the transformer more than GBDT: its paired gain is the larger one at
both larger sizes, although the intervals of the two models' gains overlap
(Table~\ref{tab:games}). With four times the games the two models have almost
the same RMSE, \V{run.gbm-hand-quadruple.score} for GBDT and
\V{run.transformer-medium-hand-quadruple.score} for the transformer.
\run{gbm-hand-quadruple} is the best model of this task, and the transformer
on four times the games does not differ significantly from the GBDT baseline
on the standard training set (gain
\V{run.transformer-medium-hand-quadruple.gain}; 95\% CI
\V{run.transformer-medium-hand-quadruple.gain.low} to
\V{run.transformer-medium-hand-quadruple.gain.high}).

The schedules of the three transformer models have the same number of
optimisation steps. On
the standard training set the transformer had its best validation RMSE at
epoch \V{run.transformer-medium-hand.best-epoch} of
\V{run.transformer-medium-hand.epochs}. On twice the games it had it at epoch
\V{run.transformer-medium-hand-double.best-epoch} of
\V{run.transformer-medium-hand-double.epochs}, and on four times in the final
epoch, \V{run.transformer-medium-hand-quadruple.best-epoch} of
\V{run.transformer-medium-hand-quadruple.epochs}, so its validation RMSE was
still decreasing when its schedule ended. We hypothesise that with four times
the data the transformer is limited by its number of optimisation steps, in
which case its gain is a lower bound. GBDT was not limited in this way: on
four times the games it used \V{run.gbm-hand-quadruple.trees} of the
\V{method.gbm.rounds} boosting rounds it was allowed.

These results are consistent with models that are limited by the number of
games. All positions of a game share one result, so the training set
contains far fewer independent outcomes than positions, and what the board
adds to the evaluation has to be learned from them.
\fi

% Feature-group ablation: the same ablation for every task, so a model paper
% inputs this file inside its ablation section and then reports its own
% findings. It appears once the task has the runs.
% The prose follows docs/writing.md.
\iffeatures
\subsection{Feature groups}
\label{sec:feature-groups}

\ifsequence
The hand-crafted features fall into three groups: move accuracy, clock use
and the number of available moves.
\else
The hand-crafted features fall into groups: material, mobility, hanging
pieces, passed pawns and king safety, and, where the task gives a line of
moves, the opponent's last move, the first move of the line and the line as a
whole.
\fi
We retrain a model with one group removed at a time
(Table~\ref{tab:features}, Figure~\ref{fig:features}). A significant loss
shows that the model used the group and that no other input provides the same
information. The absence of a loss has two possible explanations: the group
is uninformative, or the same information is available from another group or
from the other inputs. The ablation therefore measures the marginal
contribution of a group given everything else the model receives, which can
be much smaller than the information the group carries on its own.

\begin{table}[htbp]\centering\footnotesize
\block{study-features}
\caption{Feature-group ablation. Each model is trained on all hand-crafted
features except one group. The gain is paired with the model trained on all
features. A negative gain means that the model is worse without the group,
and bold marks statistical significance. Groups are ordered by the size of
the loss.}
\label{tab:features}
\end{table}

\begin{figure}[htbp]\centering
\block{study-features-figure}
\caption{Paired gain of the GBDT baseline without one feature group over the
GBDT baseline with all features, with 95\% CIs. Left of the vertical line,
the model is worse without the group.}
\label{fig:features}
\end{figure}
\fi


\iffeatures
\paragraph{Material and mobility are the groups that GBDT uses.}
The ablation was run for the GBDT baseline. It is significantly worse in
\V{study.features.gbm-hand.worse.of}. The largest loss follows the removal of
the \V{study.features.gbm-hand.most.group} group, which consists of the total
material, the material balance and the number of minor and major pieces
(paired gain \V{study.features.gbm-hand.most.gain}). Without the mobility
features, which are the numbers of legal moves, legal captures and legal
checks and whether the mover is in check, the paired gain is
\V{pair.gbm-hand-no-mobility.gbm-hand.gain}
(\V{pair.gbm-hand-no-mobility.gbm-hand.verdict}). Without the hanging-piece
features it is \V{pair.gbm-hand-no-hanging.gbm-hand.gain}
(\V{pair.gbm-hand-no-hanging.gbm-hand.verdict}), without the passed-pawn
features \V{pair.gbm-hand-no-pawns.gbm-hand.gain}
(\V{pair.gbm-hand-no-pawns.gbm-hand.verdict}), and without the two features
that count attacks near the kings \V{pair.gbm-hand-no-kings.gbm-hand.gain}
(\V{pair.gbm-hand-no-kings.gbm-hand.verdict}). Removing all hand-crafted
features gives \run{gbm-context}, whose gain relative to the GBDT baseline is
\V{run.gbm-context.gain} (Table~\ref{tab:leaderboard}).

The engine evaluation already reflects the material on the board, and yet
the material balance has the largest share of the split gain among the
hand-crafted features (\V{reference.leans.4.share}\%;
Figure~\ref{fig:leans}). We have not determined what the material features
add to the evaluation. Three explanations are possible. First, the win
probability is limited: at least one position in twenty has its maximum of
\V{describe.win-before.max}\% (Table~\ref{tab:describe}), and among such
positions the material balance may still distinguish a larger advantage from
a smaller one. Second, at the same evaluation, an advantage that consists of
extra material may be converted more reliably than one that depends on
finding particular moves. Third, the total material and the number of pieces
describe the phase of the game\iflook, and the RMSE falls with the number of
pieces on the board (Section~\ref{sec:phase})\fi. The ablation removes the
three features together and does not distinguish between these
explanations.
\fi
\fi

\section{Error Analysis}
\label{sec:error}

% Error analysis: what the analysis is and how it is computed, the same for
% every task. A model paper inputs this file at the start of its error
% analysis, before the subgroups that are specific to its task.
% The prose follows docs/writing.md.
\iflook
The main results identify the best model. They do not show where a model is
accurate and where it is not, which determines how it can be used. We
therefore analyse the test predictions of every model by subgroup. All
results in this section are computed from the \V{look.test-rows} test
examples (analysis of \V{look.date}).

\paragraph{Subgroups.}
We partition the test set by one variable at a time. For each subgroup we
report the number of test examples, the mean observed target, the mean
prediction of a model and the model's score. Two models are compared:
\V{challenger.name}, the best model other than the GBDT baseline, and the
GBDT baseline itself. The gain of the first over the second is computed
within each subgroup as for the whole test set, as a paired difference over
resampled games (\V{look.resamples} resamples), and is omitted for subgroups
with fewer than \V{look.smallest-compared} test examples. A subgroup whose
interval contains zero provides no evidence in either direction. A numeric
variable is divided into ranges, and a range written $a$-$b$ includes $a$ and
excludes $b$.

\paragraph{Two-way tables.}
For a pair of variables we show the mean observed target and the mean
prediction in each cell, on a common colour scale. Cells with fewer than
\V{grid.smallest-cell} test examples are left empty.

\paragraph{Examples.}
Example \ifsequence games\else positions\fi{} are selected by fixed rules and
not by inspection: typical examples chosen by a hash of their key, the
examples with the best model's most extreme predictions in each direction,
and its largest errors in each direction, \V{examples.count} of each. They
illustrate the quantities under discussion and are not evidence of how often
a pattern occurs.
\else
\emph{This version of the paper was built before the error analysis was run.
The analysis partitions the test set by rating, clock and properties of the
position, measures calibration, and shows example positions. It is produced
by \texttt{lab analyse} on the machine that holds the dataset and the test
predictions of every model, and this section appears once its output is
committed.}
\fi


\iflook
\paragraph{Models compared in this section.}
In this task the best model and the GBDT baseline are the same model trained
on different amounts of data, and their predictions are almost identical
(correlation \V{best.beside.correlation}). The comparison of interest is the
one between the engine-derived win probability and a fitted model. The tables
and figures of this section therefore show \run{engine-curve} next to the
GBDT baseline. Their last column, \emph{Baseline's lead}, is the RMSE of
\run{engine-curve} minus the RMSE of the GBDT baseline on the positions of
the subgroup. This difference was not resampled within subgroups and has no
confidence interval, so in a small subgroup it indicates a direction and is
not a significant result. The interval for the whole test set is given in
Section~\ref{sec:engine-baseline}. The subgroup tables for the best model,
with paired gains, are in the appendix.

\subsection{Win probability before the move}

Table~\ref{tab:win} divides the test positions by the mover's win probability
before the move, in steps of ten percentage points, and
Figure~\ref{fig:win-bias} shows the mean prediction minus the mean result for
the two models. In the four groups below 40\% the mean result is higher than
the mean win probability, and in the six groups from 40\% on it is lower. The
deviations are largest in the second group from each end. Between 10\% and
20\%, movers score \V{slice.win.10-20.happened}, where the win probability
averages \V{slice.win.10-20.engine-curve.said}, and between 80\% and 90\%
they score \V{slice.win.80-90.happened}, where it averages
\V{slice.win.80-90.engine-curve.said}. The largest deviation,
\V{slice.win.engine-curve.widest-off}, occurs in the group
\V{slice.win.engine-curve.widest-off.group}. From positions that the engine
rates between 80\% and 90\%, movers obtain a lower mean result than the
\V{slice.win.70-80.engine-curve.said} that the win probability averages over
the positions that it rates between 70\% and 80\%. In the outermost groups
the deviations are smaller: movers score \V{slice.win.-10.happened} against a
mean win probability of \V{slice.win.-10.engine-curve.said} below 10\%, and
\V{slice.win.90+.happened} against \V{slice.win.90+.engine-curve.said} at
90\% and above.

\begin{table}[htbp]\centering\footnotesize
\block{slice-win-engine-curve}
\caption{Test-set results by the mover's win probability before the move, in
percent. For each group: the number of test positions, the mean result, and
the mean prediction and RMSE of the engine-derived win probability and of the
GBDT baseline. \emph{Baseline's lead} is the difference between the two RMSE
values, without a confidence interval.}
\label{tab:win}
\end{table}

\begin{figure}[htbp]\centering
\block{slice-win-bias-engine-curve}
\caption{Mean prediction minus mean result by the mover's win probability
before the move, for the engine-derived win probability and the GBDT
baseline. Above the line a model predicts a higher score than movers
obtained.}
\label{fig:win-bias}
\end{figure}

The pattern is not symmetric about 50\%. In the two groups between 40\% and
60\% the win probability is above the mean result, by
\V{slice.win.40-50.engine-curve.off} and \V{slice.win.50-60.engine-curve.off},
and each deviation above 60\% is larger in absolute value than the one in the
mirrored group below 40\%. This asymmetry is consistent with the offset
described in Section~\ref{sec:engine-baseline}, which adds to the
overconfidence where the mover stands well and compensates part of it where
the mover stands badly.

The GBDT baseline has the lower RMSE in every group. Its lead is smallest in
the group \V{slice.win.engine-curve.behind.least.group}
(\V{slice.win.engine-curve.behind.least}) and largest in the group
\V{slice.win.engine-curve.behind.most.group}
(\V{slice.win.engine-curve.behind.most}). Its own largest deviation from the
mean result is \V{slice.win.reference.widest-off}, in the group
\V{slice.win.reference.widest-off.group}, and its deviation in the group
80-90 is \V{slice.win.80-90.reference.off}. The best model, which was trained
on \V{more.gbm-hand.most} the games, deviates by almost the same amounts in
these two groups (\V{slice.win.10-20.challenger.off} and
\V{slice.win.80-90.challenger.off}), so a larger training set does not remove
them. The two groups are related. When one player has between 10\% and 20\%,
the opponent has between 80\% and 90\%, so both groups come from games in
which one side had such an advantage, seen from the two sides, and the two
deviations are not independent. We have not determined whether they reflect
the sampling of the \V{data.games.test} test games or a property of the
models.

\subsection{Rating}

Figure~\ref{fig:grid-rating} gives the mean result by the mover's rating and
the win probability, and Figure~\ref{fig:grid-rating-off} the deviations of
the two models from it. From positions that the engine rates at 80\% or more,
movers rated below 1100 score \V{grid.rating-win.-1100.80+.happened}, where
the win probability averages
\V{grid.rating-win.-1100.80+.engine-curve.said}. The mean result rises with
every rating band, to \V{grid.rating-win.2000+.80+.happened} for movers rated
2000 and above, where the win probability averages
\V{grid.rating-win.2000+.80+.engine-curve.said}. At the other end, from
positions rated below 20\%, movers rated below 1100 score
\V{grid.rating-win.-1100.-20.happened}, where the win probability averages
\V{grid.rating-win.-1100.-20.engine-curve.said}, and the mean result falls
with every band, to \V{grid.rating-win.2000+.-20.happened} for movers rated
2000 and above (\V{grid.rating-win.2000+.-20.engine-curve.said}). From
positions rated between 60\% and 80\%, movers rated below 1100 score
\V{grid.rating-win.-1100.60-80.happened} and movers rated 2000 and above
\V{grid.rating-win.2000+.60-80.happened}.

\begin{figure}[htbp]\centering
\block{grid-rating-win}
\caption{Mean result by the mover's rating and the mover's win probability
before the move. Cells with too few test positions are left empty.}
\label{fig:grid-rating}
\end{figure}

\begin{figure}[htbp]\centering
\block{grid-rating-win-off-engine-curve}
\caption{Mean prediction minus mean result by the mover's rating and the
mover's win probability before the move. Top: the engine-derived win
probability. Bottom: the GBDT baseline. Orange cells mark a prediction above
the mean result and blue cells one below it, and both panels use the same
colour scale.}
\label{fig:grid-rating-off}
\end{figure}

The win probability is thus close to the mean result for the strongest
players and increasingly overconfident for weaker ones. This agrees with the
origin of the curve: its coefficient was fitted to games between players
rated 2300 and above~\cite{lilawinchances}, and the band of movers rated 2000
and above is the closest that this dataset comes to those players. The
largest deviation of the win probability in any cell is
\V{grid.rating-win.engine-curve.widest-off}, for movers rated
\V{grid.rating-win.engine-curve.widest-off.row} and a win probability of
\V{grid.rating-win.engine-curve.widest-off.column}, compared with
\V{grid.rating-win.gbm-hand.widest-off} for the GBDT baseline, which receives
both ratings.

These values describe games and not individual players. Most games in the
dataset are between players of similar rating
(Section~\ref{sec:rating-baseline}), so the two ends of the range show the
same games from the two sides: among lower-rated players the side with the
advantage wins less often, which lowers the result of the mover who stands
well and raises that of the mover who stands badly.

\subsection{Clock}

The engine evaluation takes no account of the clock, and the clock affects
the result. With less than 30 seconds remaining on the mover's clock, movers
score \V{slice.clock.-30.happened}, and with 30 to 60 seconds
\V{slice.clock.30-60.happened} (Table~\ref{tab:clock},
Figure~\ref{fig:clock}). The win probability averages
\V{slice.clock.-30.engine-curve.said} and
\V{slice.clock.30-60.engine-curve.said} in these two groups, so by the
evaluation these movers stand about level. With one to two minutes remaining
the mean result is \V{slice.clock.60-120.happened}, and from two minutes on
it is close to one half: \V{slice.clock.120-300.happened},
\V{slice.clock.300-480.happened} and \V{slice.clock.480+.happened} in the
three remaining groups.

\begin{table}[htbp]\centering\footnotesize
\block{slice-clock-engine-curve}
\caption{Test-set results by the time remaining on the mover's clock, in
seconds. \emph{Baseline's lead} is the difference between the two RMSE
values, without a confidence interval.}
\label{tab:clock}
\end{table}

\begin{figure}[htbp]\centering
\block{slice-clock-said-engine-curve}
\caption{Mean result and mean prediction of the engine-derived win
probability and of the GBDT baseline, by the time remaining on the mover's
clock.}
\label{fig:clock}
\end{figure}

The time remaining on the opponent's clock has the opposite effect
(Figure~\ref{fig:opponent-clock}). A mover whose opponent has less than 30
seconds remaining scores \V{slice.opponent-clock.-30.happened}, where the win
probability averages \V{slice.opponent-clock.-30.engine-curve.said}. The GBDT
baseline receives both clocks. Its mean prediction is
\V{slice.clock.-30.reference.said} for movers with less than 30 seconds and
\V{slice.opponent-clock.-30.reference.said} for movers whose opponent has
less than 30 seconds.

\begin{figure}[htbp]\centering
\block{slice-opponent-clock-said-engine-curve}
\caption{Mean result and mean prediction of the engine-derived win
probability and of the GBDT baseline, by the time remaining on the opponent's
clock.}
\label{fig:opponent-clock}
\end{figure}

The lead of the GBDT baseline over the win probability is largest under time
pressure: between \V{slice.clock.60-120.engine-curve.behind} and
\V{slice.clock.30-60.engine-curve.behind} in the three groups below two
minutes on the mover's clock, compared with
\V{slice.clock.480+.engine-curve.behind} with eight minutes or more. These
groups are small. They contain \V{slice.clock.-30.rows},
\V{slice.clock.30-60.rows} and \V{slice.clock.60-120.rows} of the
\V{look.test-rows} test positions, because most of a rapid game is played
with ample time, and the deviation of the GBDT baseline itself reaches
\V{slice.clock.reference.widest-off} in the group
\V{slice.clock.reference.widest-off.group}. The lead is positive in all
three groups, and also in the three groups below two minutes on the
opponent's clock (Table~\ref{tab:slice-opponent-clock-engine-curve} in the
appendix).

Figure~\ref{fig:grid-clock} combines the clock with the win probability.
With less than 30 seconds remaining, positions that the engine rates as
balanced, between 40\% and 60\%, yield \V{grid.clock-win.-30.40-60.happened}
for the mover (\V{grid.clock-win.-30.40-60.rows} test positions), and
positions that it rates at 80\% or more yield
\V{grid.clock-win.-30.80+.happened}, where the win probability averages
\V{grid.clock-win.-30.80+.engine-curve.said}. Ample time does not remove the
overconfidence: with five minutes or more, positions rated at 80\% or more
yield \V{grid.clock-win.300+.80+.happened}, where the win probability
averages \V{grid.clock-win.300+.80+.engine-curve.said}. The largest deviation
of the win probability in any cell, \V{grid.clock-win.engine-curve.widest-off},
occurs in a cell with \V{grid.clock-win.30-60.60-80.rows} test positions
(\V{grid.clock-win.engine-curve.widest-off.row} seconds, win probability
\V{grid.clock-win.engine-curve.widest-off.column}).

\begin{figure}[htbp]\centering
\block{grid-clock-win}

\medskip
\block{grid-clock-win-off-engine-curve}
\caption{Results by the time remaining on the mover's clock and the mover's
win probability before the move. Top: mean result. Middle and bottom: mean
prediction minus mean result for the engine-derived win probability and for
the GBDT baseline, with colours as in Figure~\ref{fig:grid-rating-off}. Cells
with too few test positions are left empty.}
\label{fig:grid-clock}
\end{figure}

\subsection{Phase of the game}
\label{sec:phase}

The RMSE of both models falls with every group of plies
(Table~\ref{tab:ply}). Before ply 20 the RMSE of the GBDT baseline is
\V{slice.ply.-20.reference} and that of the constant baseline
\V{slice.ply.-20.constant}: early in a game the result is open, and the
position, the ratings and the clocks reduce the error only slightly. From ply
80 on, the RMSE of the GBDT baseline is \V{slice.ply.80+.reference},
compared with \V{slice.ply.80+.constant} for the constant baseline. The same
pattern appears by the number of minor and major pieces on the board: the
fewer pieces remain, the lower the RMSE
(Table~\ref{tab:slice-pieces-engine-curve} in the appendix). The authors of
Allie likewise report that the outcome of a game becomes more predictable as
the game progresses, for a value head trained on human
games~\cite{allie}.

\begin{table}[htbp]\centering\footnotesize
\block{slice-ply-engine-curve}
\caption{Test-set results by ply. \emph{Baseline's lead} is the difference
between the two RMSE values, without a confidence interval.}
\label{tab:ply}
\end{table}

The lead of the GBDT baseline over the win probability grows with every
group, from \V{slice.ply.-20.engine-curve.behind} before ply 20 to
\V{slice.ply.80+.engine-curve.behind} from ply 80 on. A plausible explanation
is that late positions more often have little time on the clock and an
extreme evaluation, the two conditions under which the win probability
deviates most. The tables do not separate these factors, so this analysis may
repeat the two preceding ones.

\subsection{Hanging pieces}

The hand-crafted features include the number of hanging pieces of each side,
that is pieces attacked by a less valuable piece or attacked and not
defended. A mover with no hanging piece scores \V{slice.hanging.0.happened}
and a mover with two \V{slice.hanging.2.happened}, and a mover whose opponent
has two hanging pieces scores \V{slice.opponent-hanging.2.happened}
(Tables~\ref{tab:slice-hanging-engine-curve}
and~\ref{tab:slice-opponent-hanging-engine-curve} in the appendix). The win
probability follows these differences, because the engine's evaluation
already reflects a piece that can be captured: its means in the three groups
are \V{slice.hanging.0.engine-curve.said},
\V{slice.hanging.2.engine-curve.said} and
\V{slice.opponent-hanging.2.engine-curve.said}. The counts therefore add
little to the evaluation, in agreement with the feature-group ablation, in
which the GBDT baseline without the hanging-piece features is
\V{pair.gbm-hand-no-hanging.gbm-hand.verdict}.

\subsection{Agreement with the GBDT baseline}

Table~\ref{tab:beside} gives, for each model, the correlation of its
predictions with those of the GBDT baseline on the test set and the share of
test positions on which its prediction is closer to the result. The
correlation is \V{run.engine-curve.beside.correlation} for the engine-derived
win probability and \V{run.elo-gap.beside.correlation} for the Elo expected
score. The win probability is closer to the result than the GBDT baseline on
\V{run.engine-curve.beside.closer}\% of the positions and has the higher RMSE
nevertheless. On average it is the more extreme of the two predictions, so
it is slightly closer whenever the more likely result occurs and much further
away whenever it does not, and the squared error gives the large errors more
weight.

\begin{table}[htbp]\centering\small
\block{beside}
\caption{Agreement with the GBDT baseline on the test set: the correlation
between the two models' predictions, and the share of positions on which the
model's prediction is closer to the result than the baseline's.}
\label{tab:beside}
\end{table}

\subsection{Example positions}
\label{sec:positions}

The following examples were selected from the test set by the fixed rules
described above. Each board is shown from the mover's side. The grey arrow is
the opponent's preceding move and the blue arrow is the move that was played.
Next to each board we give the move, the result of the game for the mover,
the predictions of the best model and of the GBDT baseline, and the context
of the position, which includes the mover's win probability before the move
and the length of the game. The games are public Lichess games.

Figure~\ref{fig:typical} shows three typical positions. The predictions of
the best model lie between 0.2 and 0.52, and the results are a win, a loss
and a win. For most positions the prediction is far from every possible
result, which is why the RMSE of every model is large.

An error has a different meaning in this task than in the companion papers.
The target is the result of one game, so a model that predicts 0.9 for a
group of positions expects the mover to lose some of these games, and its
largest errors are those games. The largest overestimates
(Figure~\ref{fig:over}) are positions with a win probability of more than
75\% in games that the mover lost, and the largest underestimates
(Figure~\ref{fig:under}) are positions with a win probability below 10\% in
games that the mover won. Two of the three games in Figure~\ref{fig:over}
appear in Figure~\ref{fig:under} from the opponent's side. In one of them the
mover had more than six minutes on the clock and the opponent nine seconds,
in a game that lasted 157 plies. In the other, a player rated 2467 lost to an
opponent rated 1500. These examples show what a mean result consists of. They
do not indicate a defect of the model. The positions with the highest and the
lowest predictions are shown in the appendix (Figures~\ref{fig:highest}
and~\ref{fig:lowest}).

\examples{examples-typical}{Typical positions: the first \V{examples.count}
test positions in the order of a hash of their key.}{fig:typical}

\examples{examples-over}{The largest overestimates: the test positions with
the highest predictions of the best model among the games that the mover
lost.}{fig:over}

\examples{examples-under}{The largest underestimates: the test positions with
the lowest predictions of the best model among the games that the mover
won.}{fig:under}
\fi

\section{Practical Use}
\label{sec:use}

\begin{itemize}
\item \textbf{An expected score next to the engine's evaluation.} The
engine-derived win probability describes the position under best play. A
fitted model gives the mean result that players with these ratings and clocks
obtain from such a position. An application can display both.\iflook{} The
fitted value is accurate on average: the mean prediction of the GBDT baseline
is within \V{reference.reliability.widest-gap} of the mean result in every
bin of its predictions.\fi
\item \textbf{The difference between the two is informative.} Where the two
values differ, players of the user's rating obtain a different result from
the position than best play would: a winning position that is often not won,
or a lost position that is often saved. Such positions are candidates for a
review.
\item \textbf{GBDT is sufficient.} The GBDT baseline is trained in
\V{reference.seconds} on a CPU, and no network is significantly better.\ifmore{}
Trained on \V{more.gbm-hand.most} the games, in \V{best.seconds}, it is the
best model.\fi
\item \textbf{Most of the correction does not need the board.} GBDT on the
context features alone achieves \V{context.lift.share}\% of the GBDT
baseline's improvement over the constant baseline. Where the application
cannot run a model, a table of the mean result by rating and win probability,
with a correction for a low clock, is a possible substitute. We have not
evaluated such a table.
\item \textbf{The win probability before a move favours the mover.} An
evaluation taken before a player's move overstates that player's expected
score, on average by half the mean loss of a move
(Section~\ref{sec:engine-baseline}). A statement about how a player stood
should allow for this.
\item \textbf{Not a verdict on a single game.} No model explains more than
\V{best.r2} of the variance of the result. A prediction is the mean of many
comparable games, and the result of one game is not evidence against it.
\end{itemize}

\section{Limitations}
\label{sec:limits}

\begin{itemize}
\item \textbf{One result per game.} All positions of a game share one result,
so the precision of every score depends on the number of test games,
\V{data.games.test}. The confidence intervals allow for this by resampling
whole games, and they are wide relative to the differences between the best
models.\ifmore{} The larger training sets added training games and no test
games.\fi
\item \textbf{No confidence intervals for the win probability within
subgroups.} The difference between the engine-derived win probability and the
GBDT baseline was resampled for the whole test set. Within subgroups it is a
difference between two scores.
\item \textbf{Games decided on time.} A game lost on time counts as a loss,
whatever the position. The dataset does not record how a game ended, so the
effect of the clock on the result cannot be divided into worse moves under
time pressure and losses on time.
\item \textbf{A pairing, not a player.} The mean result of a kind of position
depends on both players. A prediction for a mover of a given rating describes
games against the opponents whom such players meet, who mostly have a similar
rating.
\item \textbf{Only analysed games.} Lichess analyses a game only when a user
requests it, so the games are a self-selected subset.
\item \textbf{The depth of the evaluations.} The evaluations are those of
Lichess's server analysis. A deeper search would change some of them.
\item \textbf{One site, one time control, one month.} The games are the first
qualifying rapid games of one month on Lichess. Ratings from another site are
on a different scale and have to be mapped first.
\item \textbf{Predictions outside the range from 0 to 1.} The fitted models
are regression models without a constraint on their output.\iflook{} The
most extreme predictions of the best model are below 0 and above 1
(Figures~\ref{fig:highest} and~\ref{fig:lowest}).\fi{} An application would
limit them to this range.
\item \textbf{Observed results, not counterfactual ones.} A model predicts
what players of a given rating score from a kind of position. It does not
predict what a particular player would have scored after a different move.
\item \textbf{Single training runs.} Each network was trained with one
random seed. The confidence intervals reflect the sampling of test games and
not the variation between training runs.
\end{itemize}

\section{Reproducibility}
\label{sec:cost}

Training the \V{runs.count} models on the leaderboard took \V{cost.seconds}
in total: \V{cost.seconds.networks} for the networks, each trained on a
single cloud GPU, and \V{cost.seconds.others} for the constant baseline, the
two baselines without fitted parameters and GBDT on a CPU. The ablation
models that are not on the leaderboard are not included in this total. The
best model consists of \V{best.trees} trees and was trained in
\V{best.seconds}. The best network, \V{network.name}, has
\V{network.parameters} parameters and was trained for \V{network.epochs.run}
epochs in \V{network.seconds} on an \V{network.device}, at about
\V{network.rows-a-second} positions per second. Table~\ref{tab:runs} lists
the hardware and training time of every model, and Table~\ref{tab:settings}
the hyperparameters of the networks of the fixed sequence.

All results are generated from committed files. The dataset is built from
its recipe by \texttt{lab data game-positions}, the models are trained by
\texttt{lab all practical-chances}, each training run writes a record to the
model's \texttt{results/} directory, and \texttt{lab paper practical-chances}
generates every number, table and figure of this paper from those records and
compiles it. A record names the build of the dataset it used
(\V{data.fingerprint}) and the commit of the code that produced it
(\V{code.commit}). Model weights and test predictions are not stored in the
repository. The model's card states where they are kept and gives the hashes
of the best model and of the GBDT baseline.

\section{Conclusion and Future Work}
\label{sec:conclusion}

The result of a rated rapid game can be predicted from a single position with
an RMSE of \V{best.score}, compared with \V{floor.score} for a constant
prediction. Most of this improvement comes from the engine evaluation of the
position. Used directly, the engine-derived win probability reaches
\V{run.engine-curve.score}, but it is overconfident, most for lower-rated
players and when little time remains, and before a move it overstates the
mover's expected score, on average by half the mean loss of a move. GBDT on
the win probability, the ratings and the clocks removes most of these
deviations.
Hand-crafted features of the position add a small further gain, and networks
that read the board add nothing that the test set can detect. Every model
leaves most of the variance of the result unexplained. The following
experiments are the natural next steps.
\begin{itemize}
\item \textbf{A recalibrated curve.} A logistic curve with one coefficient
per rating band and a term for the clock, fitted to these games, is the
smallest model that could replace the fixed curve. Its RMSE relative to
\run{engine-curve} and \run{gbm-context} would show how much of the
improvement such a curve provides.
\item \textbf{Confidence intervals within subgroups.} Resampling the
comparison between the GBDT baseline and the engine-derived win probability
within each subgroup would show which of the differences reported here are
significant.
\item \textbf{Win, draw and loss probabilities.} A model with three outputs
would show how often a balanced position is drawn at each rating, which a
single expected score does not.
\item \textbf{How the game ended.} A variable for the termination of the game
would separate losses on time from the effect of the clock on the quality of
the moves.
\item \textbf{The reliability of a rating.} The record of a game gives each
player's rating change, which is large for a new account. With this variable,
the Elo formula can be evaluated on games between players with established
ratings.
\ifmore
\item \textbf{More data together with more optimisation steps.} On
\V{more.transformer-medium-hand.most} the games the transformer had its best
validation RMSE in the final epoch. A longer schedule on that training set,
and still larger training sets for both model classes, are the missing cases.
\fi
\item \textbf{A larger test set.} The confidence intervals depend on the
number of test games. More test games are needed to distinguish the best
models from one another.
\item \textbf{A player's own games.} Before the application displays an
expected score to a player, the predictions should be compared with the
results of that player's own games, after the ratings of the player's site
have been mapped to the Lichess scale.
\end{itemize}

\bibliographystyle{unsrt}
\bibliography{../../../papers/references}

\clearpage
\appendix

\section{Hyperparameters and Features}

\begin{table}[htbp]\centering\footnotesize
\block{settings}
\caption{Hyperparameters of the networks of the fixed sequence, from their
configuration files. \emph{LR} is the peak learning rate. Values shared by
all of them are given below the table.}
\label{tab:settings}
\end{table}

\begin{table}[htbp]\centering\footnotesize
\block{hand}
\caption{Hand-crafted features: the inputs of \V{task.reference} in addition
to the scalar context features, and of a network whose configuration includes
them.}
\label{tab:hand}
\end{table}

\iflook
\section{Calibration of the Best Model}

Figure~\ref{fig:reliability-best} shows the calibration of the best model
next to that of the GBDT baseline. In the lowest of the
\V{best.reliability.groups} bins the mean prediction of the best model is
\V{best.reliability.lowest.said} and the mean result is
\V{best.reliability.lowest.happened}. In the highest bin they are
\V{best.reliability.highest.said} and \V{best.reliability.highest.happened}.
The largest difference in any bin is \V{best.reliability.widest-gap},
compared with \V{reference.reliability.widest-gap} for the GBDT baseline.

\begin{figure}[htbp]\centering
\block{reliability}
\caption{Mean result against mean prediction for the best model and the GBDT
baseline, in bins of equal size by a model's prediction. Points on the
diagonal are unbiased.}
\label{fig:reliability-best}
\end{figure}

\section{Examples with Extreme Predictions}

The three positions with the highest predictions have a win probability of
more than 94\%, and in each of them the mover is rated at least 389 points
above the opponent and won the game (Figure~\ref{fig:highest}). Two of them
come from one game, and two of the three positions with the lowest
predictions come from the same game with the other player to move. In the
third, the mover has seven seconds left against more than four minutes for
the opponent (Figure~\ref{fig:lowest}). All six predictions of the best model
lie outside the range from 0 to 1.

\examples{examples-highest}{The test positions with the highest predictions
of the best model.}{fig:highest}

\examples{examples-lowest}{The test positions with the lowest predictions of
the best model.}{fig:lowest}

\section{Subgroup Tables for the Engine-Derived Win Probability}

Each table below partitions the test positions by one variable. The columns
are the group, the number of test positions, the mean result
(\emph{Observed}), the mean prediction and the RMSE of the engine-derived win
probability and of the GBDT baseline, and the difference between the two RMSE
values (\emph{Baseline's lead}), which has no confidence interval.

\block{slices-engine-curve}

\section{Subgroup Tables for the Best Model}

The same partitions for the two models whose gain was resampled within each
subgroup: \V{challenger.name}, the best model other than the GBDT baseline,
and the GBDT baseline. The last two columns are the paired gain of the first
model over the second with its 95\% CI, in bold when statistically
significant. Of the \V{slices.compared} subgroups, the gain is significantly
positive in \V{slices.ahead} and significantly negative in
\V{slices.behind}.

\block{slices}
\fi

\end{document}
